% concatenated from RA_UPS_ML3P.tex
\documentclass[a4paper,12pt]{article}
\usepackage{amsmath,latexsym}
\usepackage{amsthm}
\usepackage{amsfonts}
\usepackage{dsfont}
\usepackage{upgreek,enumerate}
\usepackage{tipa}
\usepackage{amssymb}
\usepackage{graphicx, xcolor}
\usepackage{enumitem} 
\usepackage[square,sort,comma,numbers]{natbib}
\usepackage[utf8]{inputenc}
\usepackage{float}
\newtheorem{theorem}{Theorem}[section]
\newtheorem{corollary}[theorem] {Corollary}
\newtheorem{proposition}[theorem] {Proposition}
\newtheorem{definition}[theorem]{Definition}
\newtheorem{example}[theorem]{Example}
\newtheorem{lemma} [theorem]{Lemma}
\newtheorem{notation}[theorem]{Notation}
\theoremstyle{remark}
\newtheorem{remark}[theorem]{Remark}
\voffset=-12mm
\mathsurround=2pt
\parindent=12pt
\parskip= 4.5 pt
\lineskip=3pt
\oddsidemargin=10mm
\evensidemargin=10mm
\topmargin=55pt
\headheight=12pt
\footskip=30pt
\textheight 8.1in

\textwidth=150mm
\raggedbottom
\pagestyle{myheadings}
\hbadness = 10000
\tolerance = 10000
\let\espffile=\espfbox
\renewcommand{\refname}{
	\begin{center}
		\normalsize \bf References
	\end{center}
}
\vspace{5cm}
%\documentclass[12pt,a4paper]{article}
%\usepackage[latin1]{inputenc}
%\usepackage{amsmath}
%\usepackage{amsfonts}
%\usepackage{amssymb}
%\usepackage{graphicx}
%\usepackage[square,sort,comma,numbers]{natbib}
%\usepackage[top=0.3in, bottom=0.8in, left=0.5in, right=0.5in]{geometry}
\begin{document}
		
\begin{center}
{{\Large \textbf {Prabhakar  function and unified fractional kinetic equation in bicomplex space    % Put the title of the paper here
}}
%B	

% MESA June 06, 2026
			\medskip
			
			
			
			
{\sc  Urvashi Purohit Sharma$^{1},$\ \ Kaushik Dehingia$^{2*}$ \ \ Ritu Agarwal$^{3}$ }\\
{\footnotesize $^{1}$Department of Science and Humanities, Kalaniketan Polytechnic College Jabalpur-482001, INDIA}\\
{\footnotesize 
$^{2}$Department of Mathematics, Sonari College, Sonari 785690, Assam, India}\\
{  \footnotesize $^{3}$Department of Mathematics,
Malaviya National Institute of Technology, Jaipur-302017, INDIA}
}\\
{\footnotesize E-mail: {\it $^{1}$urvashius100@gmail.com, $^{2}$kaushikdehingia17@gmail.com, $^{3} $ragarwal.maths@mnit.ac.in}
%bicomplex Extension of the Prabhakar Function: Analyticity, Properties, transforms and Application in Fractional Kinetic Equations	

*Corresponding author
}        \end{center}
\thispagestyle{empty}
	
	
\hrulefill
\begin{abstract}
\indent
The Mittag-Leffler type functions arise naturally in the solution of fractional order integral and differential equations, especially in the investigations of the fractional generalization of the kinetic equation. This article introduces a bicomplex extension of the Prabhakar function, a generalization of the Mittag-Leffler function commonly used in fractional calculus. We explore the analyticity and determine the region of convergence for this new bicomplex Prabhakar function. Several fundamental properties are established, including its integral representations, recurrence formulas, and differential relations. Furthermore, we compute the bicomplex Laplace and Mellin transforms of the function, which are useful for solving differential and integral equations. Finally, we analyze a fractional kinetic equation where the bicomplex Prabhakar function appears both in the equation and in its solution, demonstrating its applicability in complex systems involving fractional dynamics. 
	\end{abstract}
	\hrulefill
	
	
{\small \textbf{Keywords:} }Bicomplex functions, Prabhakar function, Mittag-Leffler function, Laplace transform, Mellin transform.\\ 
\textbf{MSC 2020 Classification:} 33E12, 30G35.\\

	
\section{Introduction}

Bicomplex numbers have fascinated mathematicians for decades, offering an intriguing extension to classical complex analysis and algebra. %Originally, Cockle \cite{cockle1848, cockle1849} introduced Tessarines between 1848 and 1850, which laid the foundation for further exploration. Later, 
Segre \cite{csegre1892} formalized the notion of bicomplex numbers, extending the complex field to a multidimensional algebraic structure.

\noindent Over the years, numerous properties of bicomplex numbers have been uncovered. Researchers have delved into their algebraic, geometric, and analytical aspects, leading to significant advancements in areas such as hypercomplex analysis, differential equations, and theoretical physics \cite{catoni2008mathematics, price1991explicit, roch2014functional}. The bicomplex framework has found applications in quantum mechanics, signal processing, and applied mathematics, owing to its ability to model multidimensional wave phenomena and hyperbolic rotations \cite{ghimici2017differential, riley2001applications}.

\noindent Recent developments have further explored the functional and spectral properties of bicomplex operators, paving the way for new mathematical techniques in matrix theory, differential geometry, and control systems \cite{roch2014functional, mitelman1997bicompact, beltita2015cayley}. As the field continues to evolve, the study of bicomplex numbers remains an exciting and dynamic area, bridging the gap between classical and modern mathematical frameworks.



%Bicomplex numbers are being studied for quite a long time and a lot of work has been done in this area. Cockle \cite{cockle1848,  cockle1849} introduced \textit{Tessarines} between 1848 and 1850 following which Segre \cite{csegre1892} introduced the bicomplex numbers.\\

% points as a natural completion of the complex projective straight line. \\
\noindent %There are multiple properties of the bicomplex numbers that have been studied.  
The bicomplex numbers have been studied for a variety of properties. In recent years, researchers have tried to explore the various algebraic and geometric characteristics of bicomplex numbers and their uses.
%Many properties of the bicomplex numbers have been discovered. During the last few years researchers have aimed to study the different algebraic and geometric properties of bicomplex numbers and its applications (see,  e.g. 
\cite{  rochon2004b, ronn2001, meluna2012},
% In the recent developments,  efforts have been done to extend the integral transforms \cite{ ragarwal2017, ragarwal2017c}, %holomorphic and meromorphic functions \cite{charak2009, charak2011a}
%various theorems such as Tauberian theorem \cite{ragarwal2022bicomplex, ragarwal2015}, gap theorem \cite{ragarwal2023boundary}, a number of special functions \cite{goyal2007,goyal2006a,goyal2006}
%like  Polygamma function \cite{goyal2007}, Hurwitz Zeta function \cite{goyal2006a}, Gamma and Beta functions \cite{goyal2006},  
integral transforms \cite{ ragarwal2014a,ragarwal2017}
%  special functions \cite{goyal2007,goyal2006,rochon2004,goyal2006a}, 
in the bicomplex space from their complex counterparts. In recent research, efforts have been made to extend  the classical Mittag-Leffler function and its genralizations  in the bicomplex space \cite{agarwal2023bicomplex,ragarwal2022,ragarwal2021sept,ragarwal2023solution,bakhet2025new,bakhet2025bicomplex}. %bicomplex k-Mittag-Leffler function \cite{bakhet2025new}, bicomplex k-Mittag-Leffler functions with two parameters \cite{bakhet2025bicomplex}.



% two-parameter Mittag-Leffler function  \cite{ragarwal2021sept,ragarwal2023solution}  in the bicomplex space.\\


%Many properties of the bicomplex numbers have been discovered. During the last few years researchers have aimed to study the different algebraic and geometric properties of bicomplex numbers and its applications (see,  e.g. \cite{charak2009, rochon2004, rochon2004b, rochon2006, ronn2001, meluna2012}),
% In the recent developments,  efforts have been done to extend the integral transforms \cite{ ragarwal2017, ragarwal2017c}, %holomorphic and meromorphic functions \cite{charak2009, charak2011a}
%various theorems \cite{romesh2011,rochon2013,romesh2015a,ragarwal2022bicomplex, ragarwal2023boundary,ragarwal2015}, integral transforms \cite{ragarwal2014, ragarwal2015, ragarwal2016a, ragarwal2017, ragarwal2017c, ragarwal2022sept}  special functions \cite{goyal2007,goyal2006,rochon2004,goyal2006a}, in the bicomplex space from their complex counterparts. \\
%Magnus Gustaf Mittag–Leffler, a Swedish mathematician, developed the %Mittag–Leffer (ML) function at the start of the 19th century %\cite{mittag1905}.


%Bicomplex numbers
%is a commutative ring with unity which contains the field of complex numbers and the commutative ring of hyperbolic numbers. Bicomplex (hyperbolic) numbers are unique among the complex (real) Clifford algebras in that they are commutative but not division algebras. In fact, bicomplex numbers generalize (complexify) hyperbolic numbers \cite{rochon2006}.
%	Various generalizations and the extensions of the Mittag-Leffler function in the bicomlex space are defined by many authors %\cite{gorenflo2014,arshad2018,andric2018,roberto2018,saxena2011} which are useful in the study of fractional calculus. Mittag-Leffler function appears while studying the  fractional form of various differential equations. Recently  the bicomplex version of    Mittag-Leffler function and their  generalizations are defined by many authors 
%	\cite{ragarwal2020,ragarwal2021sept,ragarwal2022sept} 
%and fractional differential  equations in bicomplex space, bicomplex version of Mittag-Leffler function would be required.
%In this paper, efforts are made to define bicomplex version of Prabhakar function.\\
	
	
\subsection{Bicomplex Numbers}
%Segre \cite{csegre1892} defined the set of bicomplex numbers as:
The set of bicomplex numbers was defined by Segre \cite{csegre1892} as follows:
\begin{definition}[Bicomplex Number]
The set of bicomplex numbers has been defined in terms of real components as
%In terms of real components, the set of bicomplex numbers is defined as
\begin{equation}
\mathbb{T} =\{\zeta: \zeta = x_0 +i_1 x_1 +i_2 x_2 + j x_3~| ~x_0,~x_1,~x_2,~x_3~ \in \mathbb{R} \}. 	
\end{equation}
%and in terms of complex numbers it can be written as
Furthermore, it can be expressed in complex numbers as
\begin{equation}\label{eq:bc}
\mathbb{T} = \{\zeta: \zeta =w_1 + i_2w_2~ |~w_1,~w_2 ~\in \mathbb{C}\}.	
\end{equation} 
%	where $i_1$ and $i_2$ are imaginary units.
\end{definition}
\noindent The following notations will be used in this paper:
%We shall use the  notations:
$x_0=\operatorname{Re}(\zeta),~x_1=\operatorname{Im}_{i_1}(\zeta),~x_2=\operatorname{Im}_{i_2}(\zeta),~x_3=\operatorname{Im}_{j}(\zeta).$ 
	
\noindent %Segre %discussed 
%talked about the presence of zero divisors. %which he called \textit{Nullifics}.
%He noticed    that the zero divisors in bicomplex numbers constitute two ideals which he called infinite set of nullifics.
%\noindent\textbf{{Null Cone}}\\
According to \cite{rochon2004a}, the null cone is the set of all zero divisors and is defined as follows:
%The set of all zero divisors is called  null cone \cite{rochon2004a} defined as follows:
%That terminology comes from the fact that when $\zeta$ is written as $ z_1 + i_2 z_2,$ zero divisors are such that $ z_1^2 + z_2^2 = 0. $ These elements are also called singular (non-invertible) elements, otherwise it is nonsingular (invertible). The set of all singular elements of $\mathbb{T}$ is denoted by $\mathbb{NC}$ or $ \mathbb{O}_2$\\
	
\begin{equation}\label{eq:s nullcone}
\mathbb{NC} = \{w_1 +w_2i_2 ~|~w_1^2 +w_2 ^2 = 0 \}. 
\end{equation}
	
\noindent In $\mathbb{T},$ there are two non-trivial idempotent zero divisors, $e_1$ and $e_2$, which are defined as follows:
%Two non-trivial idempotent zero divisors in $\mathbb{T},$ denoted by $e_1$ and $e_2$ and defined as follows \cite{price1991}: \\

$	e_1 =  \dfrac{1 + i_1i_2}{2},~ e_2   =\dfrac  {1 - i_1i_2}{2} $, $e_1.e_2   = 0$,  $ e_1+ e_2 = 1$	   and $e_1^2=e_1,~e_2^2=e_2$.
	
\begin{definition}[Idempotent Representation]
In terms of $e_1$ and $e_2$, each element $\zeta\in \mathbb{T}$ has an unique idempotent representation which  is given by
%Every element  $\zeta\in \mathbb{T}$ has unique dempotent representation in terms of $e_1$ and $e_2$ % defined by 
\begin{equation}
\zeta = w_1 + i_2 w_2   = (w_1 - i_1 w_2) e_1 + (w_1 + i_1 w_2) e_2= \zeta_1 e_1 + \zeta_2 e_2,
\end{equation}  where  $\zeta_1 = (w_1 - i_1 w_2)$ and $\zeta_2 = (w_1 + i_1 w_2).$ 
\end{definition} 
\noindent Projection mappings  $P_1 : 	\mathbb{T} \rightarrow T_1 \subseteq \mathbb{C},$ $ P_2 : \mathbb{T} \rightarrow T_2 \subseteq \mathbb{C} $ for a bicomplex number $\zeta=z_1+i_2z_2$ are defined as  (see, e.g. \cite{rochon2004b}):\\
\begin{equation}\label{eq:p1}
P_1(\zeta)=	P_1(w_1 + i_2w_2) = P_1[(w_1 - i_1w_2)e_1 + (w_1 + i_1w_2)e_2] = (w_1 - i_1w_2) \in T_1,
\end{equation}
	%	that is
	%	\begin{equation}
	%	P_1(\zeta)=P_1(\zeta_1e_1+\zeta_2e_2)=\zeta_1,
	%	\end{equation} 
and
\begin{equation}\label{eq:p2}
P_2(\zeta)=P_2(w_1 + i_2w_2) = P_2[(w_1 -i_1w_2)e_1 + (w_1 + i_1w_2)e_2] = (w_1 + i_1w_2) \in T_2,
\end{equation}
where
\begin{equation}\label{eq:  a space A1}
T_1= \{\zeta_1= w_1-i_1w_2 \hspace{1mm}| w_1,w_2 \in \mathbb{C}\}~\text{and}~ 	T_2= \{\zeta_2= w_1+i_1w_2\hspace{1mm} | w_1,w_2 \in \mathbb{C}\}.
\end{equation}\\
Here $\mathbb{T} =T_1\times_e T_2 $ is the
product region of the bicomplex space having the components $T_1$ and $T_2.$
Set of bicomplex number contains the set of complex numbers and the set of hyperbolic numbers as a subset.
\begin{definition}[Hyperbolic Numbers]
	Hyperbolic numbers   is the set of numbers $\mathbb{D}$  (see, e.g. \cite{rochon2004b}) defined as
	\begin{equation}
	\mathbb{D} =\{x_0 +x_3j~|~x_0,~x_3 \in \mathbb{R},\, j^2=1 \}. 
		\end{equation}
\end{definition}
%	 \noindent The idempotent representation of a hyperbolic number $\vartheta=x+yj\in\mathbb{D}$ can be represented as $\vartheta=\vartheta_1e_1+\vartheta_2e_2,~\vartheta_1=x+y,~\vartheta_2=x-y.$ If $\vartheta_1,\vartheta_2\ge0$ then $\vartheta$ is said to be a nonnegative hyperbolic number, this set is represented by $\mathbb{D}^+$ and when $w_1,w_2>0$ hyperbolic number $\vartheta$ is said to be positive hyperbolic number. This concept enables the following partial order on hyperbolic numbers to be defined \cite{meluna2017}.\\
	 
\noindent	 Let     $p,q\in\mathbb{D}$ then $p\preccurlyeq q$ said that  if $q-p\in\mathbb{D}^+$ and $p \prec q$ if $q-p\in\mathbb{D}^+/\{0\}.$  
\noindent	 The role that real numbers play inside complex numbers can be compared to that of hyperbolic numbers inside bicomplex numbers. % A partial order is present in the set of hyperbolic numbers, and the bicomplex definition of the norm or modulus is hyperbolic-valued. \\

 

%	\section{Introduction}
%\noindent In the following theorem Ringleb cite{ringlab1933}, (see, e.g. \cite{riley1953})  discuss the analyticity of a bicomplex	function w.r.t. its idempotent complex component functions. This theorem plays a crucial  role while discussing the convergence of the bicomplex functions.
%\begin{theorem}[Decomposition theorem of Ringleb \cite{ringlab1933}]\label{th: ringleb}
%Let $f(\zeta)$ be analytic in a region $U\subseteq\mathbb{T},$ and let $ T_1\subseteq \mathbb{C}$ and $T_2\subseteq\mathbb{C}$ be the component regions of $\mathbb{T},$ in the $\zeta_1$ and $\zeta_2$ planes respectively. Then there exists a unique pair of complex-valued analytic functions, $f_1(\zeta_1)$ and $f_2(\zeta_2)$, defined in $ U_1\subseteq T_1$ and $U_2\subseteq T_2$, respectively, such that \\
%\begin{equation}\label{u1}
%f(\zeta) = f_1(\zeta_1)e_1 + f_2(\zeta_2) e_2,~ \zeta\in U.
%\end{equation} 
%Conversely, if $f_1(\zeta_1)$ is any complex-valued analytic function in a region $T_1$ and $ f_2(\zeta_2)$ any complex-valued analytic function in a region $T_2$ then the bicomplex-valued function $f(\zeta)$ defined by the equation (\ref{u1})  is an analytic function of the bicomplex variable $\zeta$ in the product-region $U= U_1\times_e U_2$.% (equation (\ref{T cart})).
%\end{theorem}
Price explored integration in bicomplex space in the Theorem \ref{th:int price}  \cite{price1991}. This theorem is crucial in the study of bicomplex  integrals.
%In the Theorem \ref{th:int price}, Price \cite{price1991} studied the integration in bicomplex space. This theorem plays a basic role in the study of integrals of the bicomplex function.
\begin{theorem}\label{th:int price}
	Let $U\subseteq \mathbb{T}$
	and $\Gamma_1,~\Gamma_2$ be two curves defined as 
	\begin{equation}
	\Gamma_1:w_1-i_2w_2=w_1(t)-i_1w_2(t)=\zeta_1=\zeta_1(t),\quad a\le t\le b,
	\end{equation}	
	\begin{equation}
	\Gamma_2:w_1+i_2w_2=w_1(t)+i_1w_2(t)=\zeta_2=\zeta_2(t),\quad a\le t\le b,
	\end{equation}		
	which have continuous derivatives and whose traces are in $U_1\subseteq T_1$ and $U_2\subseteq T_2$ respectively and let $\Gamma$ be the curve with trace in $U=	U_1\times_e U_2$ which is given as 
	\begin{equation}
	\Gamma:\zeta(t)=\zeta_1(t)e_1+\zeta_2(t)e_2,\quad a\le t\le b.
	\end{equation}	
	%$a\le t\le b.$
	Then the integrals of $f,~f_1,~f_2$ respectively, along the curves $\Gamma,~\Gamma_1$ and $\Gamma_2$ exist and 
	\begin{equation}
	\int_{\Gamma}f(\zeta)~d\zeta=\int_{\Gamma_1}f_1(w_1-i_1w_2)~d(w_1-i_1w_2)~e_1+\int_{\Gamma_2}f_1(w_1+i_1w_2)~d(w_1+i_1w_2)~e_2
	\end{equation}
	i.e.
	\begin{equation}
	\int_{\Gamma}f(\zeta)~d\zeta=\int_{\Gamma_1}f_1(\zeta_1)~d(\zeta_1)~e_1+\int_{\Gamma_2}f_2(\zeta_2)~d(\zeta_2)~e_2.
	\end{equation}	
\end{theorem}

\begin{definition}[Hyperbolic curve]
A (real) two-dimensional surface in $\mathbb{T}\cong_{\mathbb{R}}\mathbb{R}^4$ is called a \cite{meluna2021} hyperbolic curve $\Omega$  if its projections $\Omega e_1$ and $\Omega e_2$ on the nil-planes $\mathbb{T}{e_1}$ and $\mathbb{T}e_2,$ respectively, are usual curves, where 
$\mathbb{T}e_1=\{\zeta_1e_1:\zeta_1\in\mathbb{C}\}$ and $\mathbb{T}e_2=\{\zeta_2e_2:\zeta_2\in\mathbb{C}\}.$
\end{definition}

%\section{Preliminaries}	
%We shall also repure the following function and definitions in the main theorems.
%\noindent\textbf{} Bicomplex gamma function \\
%\noindent We would require the definition of the %bicomplex gamma function defined by  Goyal et al. \cite{goyal2006},   in the Euler product form as follows:
%\begin{equation}\label{eq gamma}
%\frac{1}{\Gamma \zeta}= \zeta e^{\gamma\zeta}\prod_{n=1}^{\infty}\left( \left( 1+\frac{\zeta}{n}\right) \exp \left( -\frac{\zeta}{n}\right) \right),\quad \zeta\in\mathbb{T},
%\end{equation}
%provided that $z_1 \neq -\left( \frac{m+l}{2}\right) ,$ and $z_2 \neq i_1\left( \frac{l-m}{2}\right) $ where $m,~l \in \mathbb{N} \cup \{0\}.$
%The Euler	constant	$\gamma(0 \leq \gamma \leq 1)$ is given by
%\begin{equation}
%\gamma= \lim\limits_{n \to \infty} (H_n-\log n),~	H_n= \sum_{k=1}^{n}\frac{1}{k}.
%\end{equation}	
%Also, in idempotent form 
%\begin{equation}
%\Gamma \zeta=\Gamma \zeta_1 e_1+\Gamma \zeta_2e_2,\quad %\zeta\in\mathbb{T},
%\end{equation}
%and in the integral form (see, e.g.\cite{goyal2006}), for $p=p_1e_1+p_2e_2,~ p_1,p_2\in\mathbb{R}^+ $
%\begin{equation}\label{eq:gamma}
%\Gamma \zeta= \int_{H}e^{-p}p^{\zeta-1}dp=\left(\int_{0}^{\infty}e^{-p_1}p_1^{\zeta_1-1}dp_1\right) e_1+\left(\int_{0}^{\infty}e^{-p_2}p_2^{\zeta_2-1}dp_2\right) e_2,
%\end{equation}
%where $H=(\gamma_1,\gamma_2)$ and $\gamma_1:0 ~\text{to}~ \infty,~\gamma_2:0 ~\text{to}~ \infty.$\\
	
\subsection{Mittag-Leffler type functions } 
In the theory of fractional calculus, Mittag-Leffler type functions are crucial. Studying Mittag-Leffler functions is primarily motivated by their significance in fractional calculus, where they serve the same important role as the exponential function in integer-order calculus.
The Mittag-Leffler function naturally occurs when fractional order differential or integral equations are solved, particularly when the fractional generalization of the kinetic equation is being studied.


%Mittag-Leffler type functions play a very important role in the theory of fractional calculus.  Studying Mittag-Leffler functions is primarily motivated by their significance in fractional calculus, where they serve the same important role as the exponential function in integer-order calculus. The Mittag-Leffler function arises naturally in the solution of fractional order integral equations or fractional order differential equations, and especially in the investigations of the fractional generalization of the kinetic equation.
Many authors have worked on different generalization of Mittag-Leffler function their properties and application of function in the area of fractional calculus. %These functions formed as a result of the development of the fractional calculus and the solution of new forms of integral and differential equations 
The developing of the fractional calculus and the finding of various integral and differential equation forms led to the development of these functions.\cite{gorenflo2014}.
        
\noindent The one parameter (classical)  Mittag-Leffler (ML) function %is a special function, It
is a direct generalization of the exponential series.  ML function is  defined by the Swedish mathematician Mittag-Leffler \cite{mittag1905,mittag1903nouvelle} at the start of the 19th century.	
%Fractional calculus theory is a generalization of differentiation and integration procedures to any non-integer order or imaginary number \cite{miller1993introduction,podlubny1998fractional}.  There are various works in the literature that are based on fractional calculus  in the context of electromagnetic theory such as fractional electromagnetic waves in conducting media \cite{gomez2016fractional}, the major goal of this study was to explain how to use the Caputo fractional derivative to solve the fractional wave equation in conducting materials, In the paper universal character of the fractional space-time electromagnetic waves in dielectric media \cite{gomez2015universal},  the fractional derivative in Caputo sense is used. Atangana-Baleanu fractional derivative applied in the work  electromagnetic waves in dielectric media \cite{gomez2016atangana}, fractional mechanical oscillators \cite{gomez2012} gives the  a simple alternative to construct fractional differential equations for physical systems. \\
	%\noindent Hence,  Mittag-Leffler function is generalization of exponential function.
	 Later, Wiman \cite{wiman1905}, defined two-parameter ML function. %is  by Wiman Agarwal and Humbert and Agarwal 	generalized and studied the properties of this function.
      It has been shown by Dzherbashyan \cite{dzherbashyan1966} that both functions are entire in $\mathbb{C}$.
    % Dzherbashyan \cite{dzherbashyan1966}, has shown that both the functions are entire in $\mathbb{C}.$  %of order $\rho=\frac{1}{\alpha},~\operatorname{Re}( \varsigma) > 0$ and type $\sigma=1.$  
	The basic properties of these functions are given in the monograph by Erdelyi et al. \cite{erdelyi1955}. %A generalized form of  Mittag-Leffler function was introduced by Prabhakar \cite{prabhakar1971}.
	Three-parameter extension of the ML function introduced by the Indian mathematician  Prabhakar \cite{prabhakar1971}. 
    This function is essential to explaining the anomalous dielectric characteristics of disordered materials, heterogeneous systems exhibiting simultaneous nonlocality and nonlinearity, and, more generally, Havriliak-Negami type models.
    %This function plays a fundamental role in the description of the 	anomalous dielectric properties in disordered materials and heterogeneous systems 	manifesting simultaneous nonlocality and nonlinearity and, more generally, in models of Havriliak-Negami type.\\	

\noindent The Prabhakar function is primarily valued for its role in describing relaxation and response in Havriliak–Negami type anomalous dielectrics, a model that accounts for the simultaneous nonlocality and nonlinearity in the behavior of disordered materials and heterogeneous systems. Additionally, the Prabhakar function finds applications in probability theory for studying stochastic processes, systems with significant anisotropy in fractional viscoelasticity, solving fractional boundary-value problems, modeling the dynamics of spherical stellar systems, and in relation to other fractional or integral differential equations. More information about the significant applications of the Prabhakar function can be found in \cite{roberto2018} and references cited therein.


\noindent The  Prabhakar function %generalized Mittag-Leffler function 
(three-parameter ML function)
is defined by Prabhakar \cite{prabhakar1971} (see e.g. \cite{gorenflo2014})
\begin{equation}\label{eq 31}
\mathbb{E}_{ \varsigma,\tau}^\delta(z)=\sum_{k=0}^{\infty}\frac{(\delta)_k}{k! \Gamma( \varsigma k+\tau )}z^k,  \hspace{2mm}z, \varsigma,\tau,\delta\in \mathbb{C},~ \operatorname{Re}( \varsigma)>0 , ~ \operatorname{Re}(\tau)>0,~ \delta\ne 0, 
\end{equation}
where $(\delta)_k= \delta (\delta+1)(\delta+2)\dots (\delta +k-1)=\dfrac{\Gamma (\delta +k)}{\Gamma \delta}.$

\noindent For $\delta=1,$ equation (\ref{eq 31}) reduces to the two-parameter  ML function  defined by \cite{wiman1905}
% \begin{equation}
% \mathbb{E}_{ \varsigma,\tau}(z)=\sum_{k=0}^{\infty}\frac{z^k}{ \Gamma( \varsigma k+\tau )},  \hspace{2mm}z\in \mathbb{C},~ \operatorname{Re}( \varsigma)>0 , ~\operatorname{ Re}(\tau)>0, ~~ \varsigma,\tau\in \mathbb{C}.
% \end{equation}	
For $\delta=1, ~ \tau =1$, equation (\ref{eq 31}) reduces to the one-parameter  ML function \cite{mittag1905}. %defined by 
% \begin{equation}\label{eq 32}
% \mathbb{E}_{ \varsigma}(z)=\sum_{k=0}^{\infty}\frac{z^k}{\Gamma( \varsigma k+1 )}  ,\hspace{2mm}z \in \mathbb{C},~ \operatorname{Re}( \varsigma)>0 ,~~  \varsigma \in \mathbb{C}.
% \end{equation}
% with $\mathbb{C}$ is the set of complex numbers.
The Prabhakar function $\mathbb{E}^\delta_{ \varsigma,\tau}(z)$ is an entire function (see, e.g. \cite[p.10]{G}) for any $ \varsigma,~\tau,~\delta\in\mathbb{C},~\operatorname{Re}( \varsigma)>0,$
and has the order $\rho$ and type $\sigma$
\begin{equation}
\rho=\frac{1}{\operatorname{Re}( \varsigma)},~\sigma=1.
\end{equation}
%(see, e.g. \cite[p.12]{saxena2011})
%\begin{equation} \varsigma\delta	E_{ \varsigma,\tau}^{\delta+1}(z)= (1+ \varsigma\delta-\tau)E_{ \varsigma,\tau}^\delta(z)+E_{ \varsigma,\tau-1}^\delta(z)\end{equation}
%Some Special Cases (see, e.g.\cite{gorenflo2014} )

\noindent In this paper, we introduce the Prabhakar function  in bicomplex space and study its properties and applications. This paper is organized as follows: Section 2, deals with definition of the Prabhakar function.  Section 3, gives the basic properties, differential relations, derivation of recurrence relations, integral relations. Section 4, deals with  Laplace transform, Mellin transform of the Prabhakar function in bicomplex space. Section 5, deals witj the the application of Prabhakar function in finding the solution of an unified fractional kinetic equation. 



\section{Bicomplex Prabhakar  Function}

Here, we introduce the bicomplex Prabhakar function   defined by 
\begin{equation}\label{eq: bc ml 3 para}
\mathbb{E}_{ \varsigma,\tau}^\delta(\zeta)=\sum_{k=0}^{\infty}\frac{(\delta)_k}{k!\Gamma( \varsigma k+\tau )}\zeta^k,  \end{equation}
		where $\zeta,  \varsigma,\tau,\delta\in \mathbb{T},$ 
$\zeta= z_1+i_2z_2= \zeta_1e_1+\zeta_2e_2$, $ \varsigma=  \varsigma_1e_1+ \varsigma_2e_2,  ~ \tau=\tau_1e_1+\tau_2e_2,~\delta=\delta_1e_1+\delta_2e_2$ with $|\operatorname{Im_j}( \varsigma)|<\operatorname{Re}( \varsigma),~|\operatorname{Im_j}(\tau)|<\operatorname{Re}(\tau).$\\


\noindent  This definition is supported by the following Theorem.
%This definition of bicomplex Prabhakar function  is well justified by the following theorem:
\begin{theorem}\label{th: ml}
Let $\zeta,  \varsigma,\tau, \delta\in \mathbb{T},$ 
$\zeta= z_1+i_2z_2= \zeta_1e_1+\zeta_2e_2$, $ \varsigma=  \varsigma_1e_1+ \varsigma_2e_2,  ~ \tau=\tau_1e_1+\tau_2e_2$ with $|\operatorname{Im_j}( \varsigma)|<\operatorname{Re}( \varsigma),~|\operatorname{Im_j}(\tau)|<\operatorname{Re}(\tau).$ 
Then
\begin{equation}
\mathbb{E}_{ \varsigma,\tau}^\delta(\zeta)=\sum_{k=0}^{\infty}\frac{(\delta)_k}{k!\Gamma( \varsigma k+\tau )}\zeta^k.  \end{equation}
		%	\begin{equation}	\begin{split}	E_{ \varsigma,\tau}(\zeta)&= \left(\sum_{k=0}^{\infty}\frac{(z_1-i_1z_2)^k}{ \delta ~(( \varsigma_1-i_1 \varsigma_2) k +( \tau_1-i_1\tau_2))} \right)e_1 \\		&+ \left(\sum_{k=0}^{\infty}\frac{(z_1+i_1z_2)^k}{\delta ~(( \varsigma_1+i_1 \varsigma_2) k +( \tau_1+i_1\tau_2))} \right)e_2. \\	&=	E_{ \varsigma,\tau}(z_1-i_1z_2)e_1+E_{ \varsigma,\tau}(z_1+i_1z_2)e_2.		\end{split}			\end{equation}
\end{theorem}
\begin{proof}
Consider the function
\begin{equation}\label{eq:ml3}
\mathbb{E}_{ \varsigma,\tau}^\delta(\zeta)=\sum_{k=0}^{\infty}\frac{(\delta)_k}{k!\Gamma( \varsigma k+\tau )}\zeta^k.  \end{equation}
By using the idempotent representation
\begin{equation}\label{eq: bc ml three para}
\begin{split}
\mathbb{E}_{ \varsigma,\tau}^\delta(\zeta)&=\sum_{k=0}^{\infty}\frac{(\delta_1)_k\zeta_1^k}{k!\Gamma( \varsigma_1 k+\tau_1 )}e_1+\sum_{k=0}^{\infty}\frac{(\delta_2)_k \zeta_2^k}{k!\Gamma( \varsigma_2 k+\tau_2 )}e_2\\
&=\mathbb{E}_{ \varsigma_1,\tau_1}^{\delta_1}(\zeta_1)e_1+\mathbb{E}_{ \varsigma_2,\tau_2}^{\delta_2}(\zeta_2)e_2,
\end{split}
\end{equation}
where $\zeta=\zeta_1 e_1+\zeta_2e_2,$  $ \varsigma= \varsigma_1e_1+ \varsigma_2 e_2,~ \tau=\tau_1e_1+\tau_2e_2$ and $\delta=\delta_1e_1+\delta_2e_2.$
		
\noindent Now,		
		
\begin{equation}
\mathbb{E}_{ \varsigma_1,\tau_1}^{\delta_1}(\zeta_1)=\sum_{k=0}^{\infty}\frac{({\delta_1})_k\zeta_1^k}{k!\Gamma( \varsigma_1 k+\tau_1 )},
\end{equation}
is the  complex Prabhakar function  convergent for $\operatorname{Re}( \varsigma_1)>0 ,~\operatorname{Re}(\tau_1)>0,~\zeta_1 \in \mathbb{C}.$\\
Similarly,
\begin{equation}\label{eq cml1}
\mathbb{E}_{ \varsigma_2,\tau_2}^{\delta_2}(\zeta_2)=\sum_{k=0}^{\infty}\frac{({\delta_2})_k \zeta_2^k}{k!\Gamma( \varsigma_2 k+\tau_2 )},
\end{equation}
is also complex Prabhakar function  convergent for $\operatorname{Re}( \varsigma_2)>0,~\operatorname{Re}(\tau_2)>0 ,~\zeta_2\in \mathbb{C}.$\\
Since $ \mathbb{E}_{ \varsigma_1,\tau_1}(\zeta_1) $ and $\mathbb{E}_{ \varsigma_2,\tau_2}(\zeta_2)$ are convergent in $T_1,~T_2$ respectively, by the  Ringleb decomposition theorem  (\ref{eq:ml3}) is also convergent in $\mathbb{T}.$\\
		
\noindent Further, Let
\begin{equation}
\begin{split}
 \varsigma&= p_0+i_1p_1+ i_2p_2+i_1i_2p_3\\
&=(p_0+i_1p_1) +i_2( p_2+i_1p_3)\\
&= \varsigma_1e_1+ \varsigma_2e_2,
\end{split}
\end{equation}
where $ \varsigma_1= (p_0+p_3) +i_1(p_1-p_2)$ and $ \varsigma_2= (p_0-p_3) +i_1(p_1+p_2).$
%and $ ~|Im( \varsigma_2)|<Re( \varsigma_1)$ 
		
%	Now$	 \varsigma=u_1+i_2u_2= \varsigma_1e_1+ \varsigma_2e_2.$\\
%Assume $u_1=x_1+i_1y_1$ and $u_2=x_2+i_1y_2$, where $x_1,~ y_1,~ x_2,~ y_2 \in \mathbb{R}.$ \\Therefore\\
\noindent Since
$\operatorname{Re}( \varsigma_1)>0$ and  $\operatorname{Re}( \varsigma_2)  >0,$
		\begin{eqnarray}
		&\Rightarrow & p_0+p_3 >0 ~\text{and}~ p_0-p_3 >0 \notag\\
		&	\Rightarrow& ~ |p_3|<p_0 \notag\\
		&\Rightarrow&|\operatorname{Im_j}( \varsigma)|<\operatorname{Re}( \varsigma).\label{eq:re(alpha)g0}
		\end{eqnarray}
		%	$\implies ~a_0+a_3 >0 $ and $a_0-a_3 >0 $\\
		%	$\implies ~ |a_3|<a_0.  $\\
		%	$\implies|\operatorname{Im_j}( \varsigma)|<\operatorname{Re}( \varsigma).$
		%	$\implies ~|Im(u_2)|<Re(u_1).$ 
		Similarly,	
		\begin{equation}
		\begin{split}
		\tau&= q_0+i_1q_1+ i_2q_2+i_1i_2q_3\\
		&=(q_0+i_1q_1) +i_2( q_2+i_1q_3)\\
		&=\tau_1e_1+\tau_2e_2,
		\end{split}
		\end{equation}
		where $\tau_1= (q_0+q_3) +i_1(q_1-q_2)$ and $\tau_2= (q_0-q_3) +i_1(q_1+q_2).$
		%and $ ~|Im( \varsigma_2)|<Re( \varsigma_1)$ 
		
		%	Now$	 \varsigma=u_1+i_2u_2= \varsigma_1e_1+ \varsigma_2e_2.$\\
		%Assume $u_1=x_1+i_1y_1$ and $u_2=x_2+i_1y_2$, where $x_1,~ y_1,~ x_2,~ y_2 \in \mathbb{R}.$ \\Therefore\\
\noindent 		Since
		$\operatorname{Re}(\tau_1)>0$ and  $\operatorname{Re}(\tau_2)  >0$,
		\begin{eqnarray}
		&\Rightarrow & q_0+q_3 >0 ~\text{and}~ q_0-q_3 >0 \notag\\
		&	\Rightarrow& ~ |q_3|<q_0 \notag \\
		&\Rightarrow&|\operatorname{Im_j}(\tau)|<\operatorname{Re}(\tau). \label{eq:re(alpha)g01}
		\end{eqnarray}
This completes the proof.		
	\end{proof}

%\noindent By substituting the value of the bicomplex gamma function defined by Goyal et al. \cite{goyal2006} in the equation (\ref{eq: bc ml three para}) we get the following representation for the Prabhakar function.

%\noindent\textbf{{Some Special Cases}}
\noindent For $\delta=1$,  \eqref{eq: bc ml 3 para} will reduce in bicomplex two parameter ML function defined by Sharma et al. \cite{ragarwal2021sept}
\begin{equation}
\mathbb{E}_{ \varsigma,\tau}^1(\zeta)=\mathbb{E}_{ \varsigma,\tau}(\zeta)=\sum_{k=0}^{\infty}\frac{\zeta^k}{\Gamma( \varsigma k+\tau )}.  \end{equation}
Also, for $\delta=1,~\tau=1$, \eqref{eq: bc ml 3 para} corresponds to the bicomplex one parameter ML function defined by Agarwal et al. \cite{ragarwal2022}
\begin{equation}
\mathbb{E}_{ \varsigma,1}^1(\zeta)=\mathbb{E}_{ \varsigma}(\zeta)=\sum_{k=0}^{\infty}\frac{\zeta^k}{\Gamma( \varsigma k+1 )}. \end{equation}



\noindent Further, by using 
the the Euler product form of bicomplex gamma function defined by  Goyal et al. \cite{goyal2006},   in   the \eqref{eq: bc ml 3 para} the bicomplex Prabhakar function can be written as:
% \begin{equation}
% \mathbb{E}_{ \varsigma,\tau}^\delta (\zeta)=\sum_{k=0}^{\infty}\frac{(\delta)_k}{k!}{\zeta^k} ( \varsigma k+\tau ) e^{\delta( \varsigma k+\tau )}\prod_{n=1}^{\infty} \left( 1+\frac{( \varsigma k+\tau )}{n}\right) exp \left( -\frac{( \varsigma k+\tau )}{n}\right).  
% \end{equation}
\begin{theorem}
	Let $\zeta,  \varsigma,\tau,\delta\in \mathbb{T}$ where $\zeta= z_1+i_2z_2= \zeta_1e_1+\zeta_2e_2$, $ \varsigma=  \varsigma_1e_1+ \varsigma_2e_2,~\tau=\tau_1e_1+\tau_2e_2$ with $|\operatorname{Im_j}( \varsigma)|<\operatorname{Re}( \varsigma),~|\operatorname{Im_j}(\tau)|<\operatorname{Re}(\tau),$ % $\delta \in \mathbb{D},~\delta\succ 0.$
	then
	\begin{equation}
	\mathbb{E}_{ \varsigma,\tau}^\delta(\zeta)=\sum_{k=0}^{\infty}\frac{(\delta)_k\zeta^k}{k!} ( \varsigma k+\tau) e^{\gamma( \varsigma k+\tau )}\prod_{n=1}^{\infty}\left( \left( 1+\frac{( \varsigma k+\tau )}{n}\right) \exp \left( -\frac{( \varsigma k+\tau )}{n}\right) \right).  
	\end{equation}
	
\end{theorem}


\begin{remark} The integral form of the bicomplex Prabhakar function is given by 
%The bicomplex Prabhakar function can be expressed in integral form using the bicomplex gamma function, which was defined by Goyal et al. \cite{goyal2006}
%Using bicomplex gamma function defined by  Goyal et al. \cite{goyal2006},   the bicomplex Prabhakar function can be represented in the integral form as
\begin{equation}
\mathbb{E}_{ \varsigma,\tau}(\zeta)= \sum_{k=0}^{\infty}\frac{(\delta)_k\zeta^k}{k!\displaystyle\int_{H}e^{-p}p^{ \varsigma k +\tau-1}dp},  
\end{equation}
%where  $H=(\gamma_1,\gamma_2)$ as defined in (\ref{eq gamma}).\\
where $H=(\gamma_1,\gamma_2)$ and $\gamma_1:0 ~\text{to}~ \infty,~\gamma_2:0 ~\text{to}~ \infty.$	
\end{remark}


\subsection{Integral Representation }
%Prabhakar function can be represented via the Mellin–Barnes integral \cite{gorenflo2014},
%If $ \varsigma\in\mathbb{R}_+,~\tau,~\delta \in\mathbb{C},~\operatorname{Re}(\tau)>0,~\operatorname{Re}(\delta)>0$ then
%\begin{equation}
%\mathbb{E}^\delta_{ \varsigma,\tau}(z)=\frac{1}{\Gamma(\gamma)}\frac{1}{2\pi i}\int_{L}\frac{\Gamma (s)\Gamma(\gamma-s)}{\Gamma(\tau- \varsigma s)}(-z)^{-s}ds,	
%\end{equation}

%\noindent  where $|\arg z|<\pi$ and the contour of integration begins at $ c-i\infty$ ends at $ c+i\infty,~0<c<\operatorname{Re}(\delta)$ and separates all poles of the integrand at $s=-k,~k=0,1,2\dots$ to the left and all poles at $s= n+\delta,~n=0,1,\dots$ to the right.\\
%Here
    We have derived the integral representation for bicomplex Prabhakar function as follows:
\begin{theorem}
    Let $ \varsigma\in \mathbb{R}_+,~\zeta,~\uplambda ,~\tau,~\delta\in \mathds{T} $ where $\zeta=\zeta_1e_1+\zeta_2e_2, ~\uplambda =\uplambda _1e_1+\uplambda _2e_2,~\tau=\tau_1e_1+\tau_2e_2 .$ Bicomplex
 Prabhakar function	can be represented as 
\begin{equation}
\mathbb{E}^\delta_{ \varsigma,\tau}(\uplambda)=\frac{1}{\Gamma(\gamma)}\frac{1}{2\pi i}\int_{\Omega}	\frac{\Gamma(\zeta)\Gamma(\delta-\zeta)}{\Gamma(\delta)\Gamma(\tau- \varsigma \zeta)}(-\uplambda )^{-\zeta}d\zeta,
\end{equation}
where	
	 $|\operatorname{Im_j}(\delta)|<\operatorname{Re}(\delta), ~|\operatorname{Im_j}(\tau)|<\operatorname{Re}(\tau)$ and $|\arg \uplambda|\prec\pi,$ hyperbolic curve $\Omega=(\Omega_1,\Omega_2).$
    The contour of integration $\Omega_r$ begins at $ c_r-i\infty$ ends at $ c_r+i\infty,~0<c_r<\operatorname{Re}(\delta_r),~(r=1,2)$ and separates all poles of the integrand at $\zeta_r
    =-p,~p=0,1,2\dots$ to the left and all poles at $\zeta_r= q+\delta,~q=0,1,\dots$ to the right (r=1,2) respectively.\\
    
%    separates all poles of the integrand at $\zeta    =-k,(k \in \mathbb{D})$ to the left plane and all poles at $\zeta= n+\delta,(n \in \mathbb{D})$ to the right plane .
    %and the contour of integration $\Omega_2$ begins at $ c-i\infty$ ends at $ c+i\infty,~0<c<\operatorname{Re}(\delta_2)$ 
   % and separates all poles of the integrand at $\zeta_2     =-k,~k=0,1,2\dots$ to the left and all poles at $\zeta_2= n+\delta,~n=0,1,\dots$ to the right.\\  \\
	
\end{theorem}
\begin{proof} By using idempotent components, we get 
	%   \begin{theorem}
	\begin{equation}\label{eq:intrep1}
	\begin{split}
&\int_{\Omega}	\frac{\Gamma(\zeta)\Gamma(\delta-\zeta)}{\Gamma(\delta)\Gamma(\tau- \varsigma \zeta)}(-\uplambda )^{-\zeta}d\zeta\\
&=\int_{\Omega_1}	\frac{\Gamma(\zeta_1)\Gamma(\delta_1-\zeta_1)}{\Gamma(\delta_1)\Gamma(\tau_1- \varsigma_1 \zeta_1)}(-\uplambda_1 )^{-\zeta_1}d\zeta_1~e_1+\int_{\Omega_2}	\frac{\Gamma(\zeta_2)\Gamma(\delta_2-\zeta_2)}{\Gamma(\delta_2)\Gamma(\tau_2- \varsigma_2 \zeta_2)}(-\uplambda_2 )^{-\zeta_2}d\zeta_2~ e_2,
	\end{split}        
	\end{equation}
    where $|\arg \uplambda_1|<\pi,~|\arg \uplambda_2 |<\pi.$ \\ 
%    The contour of integration $\Omega_1$ begins at $ c_1-i\infty$ ends at $ c_1+i\infty,~0<c_1<\operatorname{Re}(\delta_1)$ and separates all poles of the integrand at $\zeta_1     =-k,~k=0,1,2\dots$ to the left and all poles at $\zeta_1= n+\delta_1,~n=0,1,\dots$ to the right.
    
 % \noindent    The contour of integration $\Omega_2$ begins at $ c_2-i\infty$ ends at $ c_2+i\infty,~0<c<\operatorname{Re}(\delta_2)$ and separates all poles of the integrand at $\zeta_2     =-k,~k=0,1,2\dots$ to the left and all poles at $\zeta_2= n+\delta_2,~n=0,1,\dots$ to the right. \\
    \noindent Luna-Elizarrar\'{a}s et al. \cite{meluna2021} discussed residue theorem in the bicomplex space. Here  $\Omega=(\Omega_1,\Omega_2)$ is hyperbolic curve.
\noindent	Since the value of the integral $\displaystyle\dfrac{1}{2\pi i}\int_{\Omega_1}	\dfrac{\Gamma(\zeta_1)\Gamma(\delta_1-\zeta_1)}{\Gamma(\delta_1)\Gamma(\tau_1- \varsigma_1 \zeta_1)}(-\uplambda_1 )^{-\zeta_1}d\zeta_1$ is equal to   the sum of residues at the poles \cite{gorenflo2014}
	$\zeta_1=0,-1,-2,\dots$ in the complex space \\
	\begin{equation}\label{eq:intrep2}
	\begin{split}
	\frac{1}{2\pi i}\int_{\Omega_1}	\frac{\Gamma(\zeta_1)\Gamma(\delta_1-\zeta_1)}{\Gamma(\delta_1)\Gamma(\tau_1- \varsigma_1 \zeta_1)}(-\uplambda_1 )^{-\zeta_1}d\zeta_1&=\sum_{k=0}^{\infty}\lim_{\zeta_1\to k}\frac{(\zeta_1+k)\Gamma(\zeta_1)\Gamma(\delta_1-\zeta_1)(-\uplambda_1 )^{-\zeta_1}}{\Gamma(\tau_1- \varsigma_1 \zeta_1)}\\
	&=\sum_{k=0}^{\infty}\frac{(-1)^k}{k !}\frac{\Gamma(\delta_1+k)}{\Gamma(\tau_1+ \varsigma_1k)} (-\uplambda_1 )^{k}\\
		&=\Gamma(\delta_1)\sum_{k=0}^{\infty}\frac{(\delta_1)_k}{\Gamma(\tau_1+ \varsigma_1k)}\frac{ (\uplambda_1 )^{k}}{k!}\\
		&=\Gamma(\delta_1)	\mathbb{E}^{\delta_1}_{ \varsigma_1,\tau_1}(\uplambda_1).
		\end{split}
	\end{equation}
	Similarly, 
	\begin{equation}\label{eq:intrep3}
\frac{1}{2\pi i}\int_{\Omega_2}	\frac{\Gamma(\zeta_2)\Gamma(\delta_2-\zeta_2)}{\Gamma(\delta_2)\Gamma(\tau_2- \varsigma_2 \zeta_2)}(-\uplambda_2 )^{-\zeta_2}d\zeta_2=\Gamma(\delta_2)	\mathbb{E}^{\delta_2}_{ \varsigma_2,\tau_2}(\uplambda_2).
	\end{equation}
	By using the equations (\ref{eq:intrep2}) and (\ref{eq:intrep3}) in equation (\ref{eq:intrep1}) we get 
		\begin{equation}
	\begin{split}
	\frac{1}{2\pi i}\int_{\Omega}	\frac{\Gamma(\zeta)\Gamma(\delta-\zeta)}{\Gamma(\delta)\Gamma(\tau- \varsigma \zeta)}(-\uplambda )^{-\zeta}d\zeta&=\Gamma(\delta_1)	\mathbb{E}^{\delta_1}_{ \varsigma_1,\tau_1}(\uplambda_1) e_1+\Gamma(\delta_2)	\mathbb{E}^{\delta_2}_{ \varsigma_2,\tau_2}(\uplambda_2) e_2\\
	&=\Gamma(\delta)	\mathbb{E}^\delta_{ \varsigma,\tau}(\uplambda).
	\end{split}        
	\end{equation}
	
\end{proof}

\begin{theorem}
	The bicomplex Prabhakar function satisfies bicomplex Cauchy-Riemann equation.
\end{theorem}
\begin{proof}
	%By the result (\ref{eq: bc ml three para}) we have
    According to (\ref{eq: bc ml three para}), we have
	\begin{equation}
	\begin{split}
	\mathbb{E}^\delta_{ \varsigma,\tau}(\zeta)&=\mathbb{E}^{\delta_1}_{ \varsigma_1,\tau_1}(\zeta_1)e_1+\mathbb{E}^{\delta_2}_{ \varsigma_2,\tau_2}(\zeta_2)e_2\\
	&=\mathbb{E}^{\delta_1}_{ \varsigma_1,\tau_1}(z_1 - i_1z_2)e_1+\mathbb{E}^{\delta_2}_{ \varsigma_2,\tau_2}(z_1 + i_1z_2)e_2\\
	&=\mathbb{E}^{\delta_1}_{ \varsigma_1,\tau_1}(z_1 - i_1z_2)\left( \frac{1 +i_1i_2}{2}\right) +\mathbb{E}^{\delta_2}_{ \varsigma_2,\tau_2}(z_1 + i_1z_2)\left( \frac{1 - i_1i_2}{2}\right) \\
	&=\left( \frac{1}{2}\left( \mathbb{E}^{\delta_1}_{ \varsigma_1,\tau_1}(z_1 - i_1z_2)+\mathbb{E}^{\delta_2}_{ \varsigma_2,\tau_2}(z_1 +  i_1z_2)\right) \right)\\
	&~~~+i_2 \left( \frac{i_1}{2}\left( \mathbb{E}^{\delta_1}_{ \varsigma_1,\tau_1}(z_1 - i_1z_2)-\mathbb{E}^{\delta_2}_{ \varsigma_2,\tau_2}(z_1 +  i_1z_2)\right) \right)\\
	&=f_1(z_1,z_2)+i_2f_2(z_1,z_2).\\
	\end{split}
	\end{equation}
	where\\
	$f_1(z_1,z_2)=\left( \frac{1}{2}\left( \mathbb{E}^{\delta_1}_{ \varsigma_1,\tau_1}(z_1 - i_1z_2)+\mathbb{E}^{\delta_2}_{ \varsigma_2,\tau_2}(z_1 +  i_1z_2)\right) \right),$\\
	$f_2(z_1,z_2)= \left( \frac{i_1}{2}\left( \mathbb{E}^{\delta_1}_{ \varsigma_1,\tau_1}(z_1 - i_1z_2)-\mathbb{E}^{\delta_2}_{ \varsigma_2,\tau_2}(z_1 +  i_1z_2)\right) \right).$\\
	\begin{eqnarray}
	\frac{\partial f_1}{\partial z_1 }&=&\frac{1}{2}\left( \left( \mathbb{E} _{ \varsigma_1,\tau_1}^{\delta_1 }(z_1 - i_1z_2)\right) '+\left( \mathbb{E}^{\delta_2}_{ \varsigma_2,\tau_2}(z_1 +  i_1z_2)\right) '\right), \nonumber\\
	\frac{\partial f_1}{\partial z_2 }&=&\frac{-i_1}{2}\left(\left(  \mathbb{E}^{\delta_1}_{ \varsigma_1,\tau_1}(z_1 - i_1z_2)\right) '-\left( \mathbb{E}^{\delta_2}_{ \varsigma_2,\tau_2}(z_1 +  i_1z_2)\right)'\right) ,\nonumber\\
	\frac{\partial f_2}{\partial z_1 }&=&\frac{i_1}{2}\left(\left(  \mathbb{E}^{\delta_1}_{ \varsigma_1,\tau_1}(z_1 - i_1z_2)\right) '-\left( \mathbb{E}^{\delta_2}_{ \varsigma_2,\tau_2}(z_1 +  i_1z_2)\right)'\right) ,\nonumber\\
	\frac{\partial f_2}{\partial z_2 }&=&\frac{1}{2}\left(\left(  \mathbb{E}^{\delta_1}_{ \varsigma_1,\tau_1}(z_1 - i_1z_2)\right) '+\left( \mathbb{E}^{\delta_2}_{ \varsigma_2,\tau_2}(z_1 +  i_1z_2)\right)'\right) .\nonumber 
	\end{eqnarray}
	from the above equations it can be shown that
	\begin{equation}
	\frac{\partial f_1}{\partial z_1} = \frac{\partial f_2}{\partial z_2} \hspace{5mm}\text{and} \hspace{5mm}
	\frac{\partial f_2}{\partial z_1} = - \frac{\partial f_1}{\partial z_2}.
	\end{equation}
%	Hence bicomplex Cauchy-Riemann equations are satisfied by the bicomplex Prabhakar function.
Hence bicomplex Cauchy-Riemann equations are satisfied.
\end{proof}
\section{Properties of bicomplex Prabhakar function}
%Some special cases 

%\item 
%	$\phi( \varsigma,\tau,\zeta)=\delta (\tau)E_{1,\tau}^\delta(\zeta).$		where $\phi( \varsigma,\tau,\zeta)$ is the Kummer?s confluent hypergeometric function.

\noindent\textbf{Recurrence Relation}\\
The Prabhakar function satisfies the following recurrence relation  (see, e.g.\cite[p.272]{G})
\begin{equation}\label{eq: p rec}
z\mathbb{E}^\delta_{ \varsigma, \varsigma+\tau}(z)=\mathbb{E}^\delta_{ \varsigma,\tau}(z)-\mathbb{E}^{\delta-1}_{ \varsigma,\tau}(z),~\operatorname{Re}(\tau)>0,
\end{equation}


\begin{equation}\label{eq:p3}
(\tau- \varsigma\delta-1)\mathbb{E}^{\delta}_{ \varsigma,\tau}(z)=\mathbb{E}^{\delta}_{ \varsigma,\tau-1}(z)- \varsigma\delta\mathbb{E}^{\delta+1}_{ \varsigma,\tau}(z),~\operatorname{Re}(\tau)>1,
\end{equation}

\noindent where  $z, \varsigma,\tau\in\mathbb{C}, ~\operatorname{Re}( \varsigma)>0.$\\
Extension of these recurrence relations \eqref{eq: p rec} and \eqref{eq:p3} in bicomplex space is given by following theorem:
\begin{theorem}[Recurrence Relation]\label{th: bc rec}
	Let $\zeta,~ \varsigma,~\tau \in\mathbb{T}$ then for $|\operatorname{Im_j}( \varsigma)|<\operatorname{Re}( \varsigma),$
    \begin{enumerate}[label=(\roman*)]
    \item $\zeta\mathbb{E}^\delta_{ \varsigma, \varsigma+\tau}(\zeta)=\mathbb{E}^\delta_{ \varsigma,\tau}(\zeta)-\mathbb{E}^{\delta-1}_{ \varsigma,\tau}(\zeta),~|\operatorname{Im_j}(\tau)|<\operatorname{Re}(\tau).$
        \item  
        $(\tau- \varsigma\delta-1)\mathbb{E}^{\delta}_{ \varsigma,\tau}(\zeta)=\mathbb{E}^{\delta}_{ \varsigma,\tau-1}(\zeta)- \varsigma\delta\mathbb{E}^{\delta+1}_{ \varsigma,\tau}(\zeta),~|\operatorname{Im_j}(\tau)|<\operatorname{Re}(\tau)-1.
$  \end{enumerate}


\end{theorem}
\begin{proof}
Writing $\zeta=\zeta_1 e_1+\zeta_2e_2,$  $ \varsigma= \varsigma_1e_1+ \varsigma_2 e_2$ and $ \tau=\tau_1e_1+\tau_2e_2.$
From the equations (\ref{eq: bc ml three para}) and (\ref{eq: p rec}), we have 
	\begin{equation}
	\begin{split}
		\zeta\mathbb{E}^\delta_{ \varsigma, \varsigma+\tau}(\zeta)&=\left( \zeta_1\mathbb{E}^\delta_{ \varsigma_1, \varsigma_1+\tau_1}(\zeta_1)\right) e_1+\left( \zeta_2\mathbb{E}^\delta_{ \varsigma_2, \varsigma_2+\tau_2}(\zeta_2)\right) e_2\\
	&=\left( \mathbb{E}^\delta_{ \varsigma_1,\tau_1}(\zeta_1)-\mathbb{E}^{\delta-1}_{ \varsigma_1,\tau_1}(\zeta_1)\right)e_1 +\left( \mathbb{E}^\delta_{ \varsigma_2,\tau_2}(\zeta_2)-\mathbb{E}^{\delta-1}_{ \varsigma_2,\tau_2}(\zeta_2)\right)e_2\\
	&=\mathbb{E}^\delta_{ \varsigma_1e_1+ \varsigma_2e_2,\tau_1e_1+\tau_2e_2}(\zeta_1e_1+\zeta_2e_2)-\mathbb{E}^{\delta-1}_{ \varsigma_1e_1+ \varsigma_2e_2,\tau_1e_1+\tau_2e_2}(\zeta_1e_1+\zeta_2e_2)\\
		&=\mathbb{E}^\delta_{ \varsigma,\tau}(\zeta)-\mathbb{E}^{\delta-1}_{ \varsigma,\tau}(\zeta).
		\end{split}
	\end{equation}

\noindent (ii) From the results (\ref{eq:p3}) and (\ref{eq: bc ml three para}) we have 
	\begin{equation}
	\begin{split}
	(\tau- \varsigma\delta-1)\mathbb{E}^{\delta}_{ \varsigma,\tau}(\zeta)&=(\tau_1- \varsigma_1\delta-1)\mathbb{E}^{\delta}_{ \varsigma_1,\tau_1}(\zeta_1)e_1+(\tau_2- \varsigma_2\delta-1)\mathbb{E}^{\delta}_{ \varsigma_2,\tau_2}(\zeta_2)e_2\\
&=\left(\mathbb{E}^{\delta}_{ \varsigma_1,\tau_1-1}(\zeta_1)- \varsigma_1\delta\mathbb{E}^{\delta+1}_{ \varsigma_1,\tau_1}(\zeta_1)\right) e_1	\\
&~~~+\left(\mathbb{E}^{\delta}_{ \varsigma_2,\tau_2-1}(\zeta_2)- \varsigma_2\delta\mathbb{E}^{\delta+1}_{ \varsigma_2,\tau_2}(\zeta_1)\right) e_2	\\
&=\mathbb{E}^{\delta}_{ \varsigma_1e_1+ \varsigma_2e_2,\tau_1e_1+\tau_2e_2-1}(\zeta_1e_1+\zeta_2e_2)\\
&~~~-( \varsigma_1e_1+ \varsigma_2e_2)\delta\mathbb{E}^{\delta+1}_{ \varsigma_1e_1+ \varsigma_2e_2,\tau_1e_1+\tau_2e_2}(\zeta_1e_1+\zeta_2e_2)	\\
&=\mathbb{E}^{\delta}_{ \varsigma,\tau-1}(\zeta)- \varsigma\delta\mathbb{E}^{\delta+1}_{ \varsigma,\tau}(\zeta).
\end{split}
\end{equation}
Here,
\begin{equation}
		\begin{split}
		\tau&= q_0+i_1q_1+ i_2q_2+i_1i_2q_3\\
		&=(q_0+i_1q_1) +i_2( q_2+i_1q_3)\\
		&=\tau_1e_1+\tau_2e_2,
		\end{split}
		\end{equation}
		where $\tau_1= (q_0+q_3) +i_1(q_1-q_2)$ and $\tau_2= (q_0-q_3) +i_1(q_1+q_2).$
		
\noindent 		Since
		$\operatorname{Re}(\tau_1)>1$ and  $\operatorname{Re}(\tau_2)  >1$
		\begin{eqnarray}\notag
		&\Rightarrow & q_0+q_3 >1 ~\text{and}~ q_0-q_3 >1.\\\notag
        &\Rightarrow & q_3 >1-q_0 ~\text{and}~ -q_3 >1-q_0.\\\notag
 &\Rightarrow & -q_3<q_0-1 ~\text{and}~ q_3 <q_0-1.\\\notag
		&	\Rightarrow& ~ |q_3|<q_0-1.\label{eq:re(alpha)g02}\\
		&\Rightarrow&|\operatorname{Im_j}(\tau)|<\operatorname{Re}(\tau)-1.
		\end{eqnarray}

\end{proof}




\noindent\textbf{Differential recurrence relations} \\
Differential recurrence relations for the complex Prabhakar function are  (see, e.g.\cite[p.91]{mathai2008})
\begin{equation}\label{eq:p11}
\left( \frac{d}{dz}\right)^n\mathbb{E}^{\delta}_{ \varsigma,\tau}(z)= (\delta)_n \mathbb{E}^{\delta+n}_{ \varsigma,\tau+n \varsigma}(z), 
\end{equation}
and
\begin{equation}\label{eq:p21}
\left( z\frac{d}{dz}+\delta\right)\mathbb{E}^{\delta}_{ \varsigma,\tau}(z)= \delta \mathbb{E}^{\delta+1}_{ \varsigma,\tau}(z),  \end{equation}
where $ ~\operatorname{Re}( \varsigma)>0,~\operatorname{Re}(\tau)>0.$


\noindent Another significant differential relation of the complex Prabhakar function is given by (see, e.g.\cite[p.99]{gorenflo2014})
\begin{equation}\label{eq: p31}
\left( \frac{d}{dz}\right)^n\left( z^{\tau-1}\mathbb{E}^\delta_{ \varsigma,\tau} (\omega z^ \varsigma) \right) =z^{\tau-n-1}\mathbb{E}^\delta_{ \varsigma,\tau-n}(\omega z^ \varsigma),~n=1,2,3 \dots,
\end{equation}
where $ \varsigma,\tau,\delta,z,\omega \in\mathbb{C},~\operatorname{Re}(\tau)>n.$ 

% In particular
% \begin{equation}
% \frac{d^n}{dz^n}\left( z^{\tau-1}\mathbb{E}_{ \varsigma,\tau}(az^ \varsigma)\right) =z^{\tau-n-1}\mathbb{E}_{ \varsigma,\tau-n}(az^ \varsigma).
% \end{equation}

%\noindent In the following theorem, we develop the differential recurrence relations for the bicomplex Prabhakar function.
\noindent Similar to Theorem  \ref{th: bc rec}, breaking in the idempotent components and using the  above results  \eqref{eq:p11}, \eqref{eq:p21} and \eqref{eq: p31} we obtain following differential recurrence relations for the bicomplex Prabhakar function 
 

\begin{theorem}
	Let $\zeta \in\mathbb{T}$, $~|\operatorname{Im_j}(\tau)|<\operatorname{Re}(\tau)$, $|\operatorname{Im_j}( \varsigma)|<\operatorname{Re}( \varsigma)$ then bicomplex Prabhakar functions satisfies the following  differential recurrence relations:
	\begin{enumerate}[label=(\roman*)]
		\item 
		$\left( \frac{d}{d\zeta}\right)^n\mathbb{E}^{\delta}_{ \varsigma,\tau}(\zeta)= (\delta)_n \mathbb{E}^{\delta+n}_{ \varsigma,\tau+n \varsigma}(\zeta).$ 
		\item 
		$\left( \zeta\frac{d}{d\zeta}+\delta\right)\mathbb{E}^{\delta}_{ \varsigma,\tau}(\zeta)= \delta \mathbb{E}^{\delta+1}_{ \varsigma,\tau}(\zeta). $
\item If $ \varsigma,~\tau,~\delta,~\zeta,~\uplambda \in\mathbb{T}$ then for $|\operatorname{Im_j}(\tau)|<\operatorname{Re}(\tau)-n$
\begin{equation}
\left( \frac{d}{d\zeta}\right)^n\left( \zeta^{\tau-1}\mathbb{E}^\delta_{ \varsigma,\tau} (\uplambda \zeta^ \varsigma) \right) =\zeta^{\tau-n-1}\mathbb{E}^\delta_{ \varsigma,\tau-n}(\uplambda \zeta^ \varsigma),~n=1,2,3\dots.
\end{equation}
	\end{enumerate}
    
\end{theorem}
%\begin{proof}
%\noindent (i)
%From the results (\ref{eq:p1}) and (\ref{eq: bc ml three para}), we have for $|\operatorname{Im_j}( \varsigma)|<\operatorname{Re}( \varsigma),~|\operatorname{Im_j}(\tau)|<\operatorname{Re}(\tau)$
%	\begin{equation}
%	\begin{split}
%	\left( \frac{d}{d\zeta}\right)^n\mathbb{E}^{\delta}_{ \varsigma,\tau}(\zeta)&=\left( \frac{d}{d\zeta_1}\right)^n\mathbb{E}^{\delta}_{ \varsigma_1,\tau_1}(\zeta_1)e_1+\left( \frac{d}{d\zeta_2}\right)^n\mathbb{E}^{\delta}_{ \varsigma_2,\tau_2}(\zeta_2)e_2\\
%	&=(\delta)_n \mathbb{E}^{\delta+n}_{ \varsigma_!,\tau_1+n \varsigma_1}(\zeta_1)e_1+(\delta)_n \mathbb{E}^{\delta+n}_{ \varsigma_2,\tau_2+n \varsigma_2}(\zeta_2)e_2\\
	%&=(\delta)_n \mathbb{E}^{\delta+n}_{ \varsigma_1e_1+ \varsigma_2e_2,\tau_1e_1+\tau_2e_2+n( \varsigma_1e_1+ \varsigma_2e_2)}(\zeta_1e_1+\zeta_2e_2)\\
	%&=(\delta)_n \mathbb{E}^{\delta+n}_{ \varsigma,\tau+n \varsigma}(\zeta).
	%	\end{split}
	%\end{equation}
	
%\noindent (ii) From the results (\ref{eq:p2}) and (\ref{eq: bc ml three para}) we have for $|\operatorname{Im_j}( \varsigma)|<\operatorname{Re}( \varsigma),~|\operatorname{Im_j}(\tau)|<\operatorname{Re}(\tau)$
%	\begin{equation}
%	\begin{split}
	%\left(\zeta\frac{d}{d\zeta}+\delta\right)\mathbb{E}^{\delta}_{ \varsigma,\tau}(\zeta)&=\left(\zeta_1\frac{d}{d\zeta_2}+\delta\right)\mathbb{E}^{\delta}_{ \varsigma_1,\tau_1}(\zeta_1)e_1+\left(\zeta_2\frac{d}{d\zeta_2}+\delta\right)\mathbb{E}^{\delta}_{ \varsigma_2,\tau_2}(\zeta_2)e_2\\
	%&= \delta \mathbb{E}^{\delta+1}_{ \varsigma_1,\tau_1}(\zeta_1)e_1+\delta \mathbb{E}^{\delta+1}_{ \varsigma_2,\tau_2}(\zeta_2)e_2\\
%&=	\delta \mathbb{E}^{\delta+1}_{ \varsigma_1e_1+ \varsigma_2e_2,\tau_1e_1+\tau_2e_2}(\zeta_1e_1+\zeta_2e_2)\\
%&=\delta \mathbb{E}^{\delta+1}_{ \varsigma,\tau}(\zeta).
%	\end{split}
%	\end{equation}

%\noindent (iii)
%	\begin{equation}
%	\begin{split}
%	\left( \frac{d}{d\zeta}\right)^n\left( \zeta^{\tau-1}\mathbb{E}^\delta_{ \varsigma,\tau} (\omega \zeta^ \varsigma) \right)&=
%	\left( \frac{d}{d\zeta_1}\right)^n\left( \zeta_1^{\tau_1-1}\mathbb{E}^\delta_{ \varsigma_1,\tau_1} (\omega_1 \zeta_1^{ \varsigma_1}) \right)e_1 +	\left( \frac{d}{d\zeta_2}\right)^n\left( \zeta_2^{\tau_2-1}\mathbb{E}^\delta_{ \varsigma_2,\tau_2} (\omega_2 \zeta_2^{ \varsigma_2}) \right)e_2\\	
%	 &=\zeta_1^{\tau_1-n-1}\mathbb{E}^\delta_{ \varsigma_1,\tau_1-n}(\omega_1 \zeta_1^{ \varsigma_1})e_1+\zeta_2^{\tau_2-n-1}\mathbb{E}^\delta_{ \varsigma_2,\tau_2-n}(\omega_2 \zeta_2^{ \varsigma_2})e_2\\
%	 &=(\zeta_1e_+\zeta_2e_2)^{\tau_1e_1+\tau_2e_2-n-1}\\
	 %&~~~~~\mathbb{E}^\delta_{ \varsigma_1e_1+ \varsigma_2e_2,\tau_1e_1+\tau_2e_2-n}((\omega_1e_1+\omega_2e_2)( \zeta_1e_1
	 %+\zeta_2e_2))^{ \varsigma_1e_1+ \varsigma_2e_2}\\
%&=\zeta^{\tau-n-1}\mathbb{E}^\delta_{ \varsigma,\tau-n}(\omega \zeta^ \varsigma).
%\end{split}
%\end{equation}
%	Also
\begin{remark}
   Since  $\operatorname{Re}(\tau_1)>n$ and  $\operatorname{Re}(\tau_2)  >n$
	\begin{eqnarray}\notag
	&\Rightarrow & q_0+q_3 >n ~\text{and}~ q_0-q_3 >n.\notag\\
	&\Rightarrow & q_3 >n-q_0 ~\text{and}~ -q_3 >n-q_0.\notag\\
		&\Rightarrow & -q_3 <q_0-n ~\text{and}~ q_3 <q_0-n.\notag\\
	&	\Rightarrow& ~ |q_3|<q_0-n.\label{eq:re(alpha)g}\notag\\
	&\Rightarrow&|\operatorname{Im_j}(\tau)|<\operatorname{Re}(\tau)-n.
	\end{eqnarray}
	
\end{remark}	
%\end{proof}

\noindent\textbf{Integral Relation}\\
Corresponding to the integral relation for the complex Prabhakar  function (see, e.g.\cite[p.100]{gorenflo2014})
\begin{equation}\label{eq:p int}
\int_{0}^{z}t^{\tau-1}\mathbb{E}^\delta_{ \varsigma,\tau}(\omega t^ \varsigma)dt=z^{\tau}\mathbb{E}^\delta_{ \varsigma,\tau+1}(\omega z^ \varsigma),
\end{equation}
where $ \varsigma,~\tau,~\delta,~z,~\omega\in\mathbb{C},~\operatorname{Re}( \varsigma)>0,~\operatorname{Re}(\tau)>0,~\operatorname{Re}(\delta)>0,$ we obtain the integral relation for the bicomplex Prabhakar function contained in the following theorem.\\
% In particular
% \begin{equation}
% \int_{0}^{z}t^{\tau-1}\mathbb{E}_{ \varsigma,\tau}(ut^ \varsigma)dt=z^{\tau}\mathbb{E}_{ \varsigma,\tau+1}(uz^ \varsigma).
% \end{equation}

\noindent Breaking the integral into idempotent components and using the result \eqref{eq:p int}, similar to  Theorem \ref{th: bc rec}, we get the following  relation. % for the bicomplex Prabhakar function.
\begin{theorem}[Integral Relation]
	If $ \varsigma,~\tau,~\delta,~\zeta,~\uplambda \in\mathbb{T}$ with $|\operatorname{Im_j}( \varsigma)|<\operatorname{Re}( \varsigma),~|\operatorname{Im_j}(\tau)|<\operatorname{Re}(\tau),~|\operatorname{Im_j}(\delta)|<\operatorname{Re}(\delta)$ then
	
	%\begin{equation}\label{eq:p int1}
%	\int_{0}^{\zeta}t^{\tau-%1}\mathbb{E}^\delta_{ \varsigma,\tau}(\omega %t^ \varsigma)dt=\zeta^{\tau}\mathbb{E}^\delta_{\upal%pha,\tau+1}(\omega \zeta^ \varsigma).
%	\end{equation}
    \begin{equation}\label{eq:p int1}
	\int_{\Gamma}t^{\tau-1}\mathbb{E}^\delta_{ \varsigma,\tau}(\uplambda t^ \varsigma)dt=\zeta^{\tau}\mathbb{E}^\delta_{ \varsigma,\tau+1}(\uplambda  \zeta^ \varsigma),
    \end{equation}
where $\Gamma =(\Gamma_1.\Gamma_2) ,~\Gamma_1:0 ~\text{to}~ \zeta_1,~\Gamma_2:0 ~\text{to}~ \zeta_2.$ 
%is piecewise continuous differentiable closed
% contour in the bicomplex space	
\end{theorem}

%\begin{proof} From the results (\ref{eq:p int}) and (\ref{eq: bc ml three para}) we have for $|\operatorname{Im_j}( \varsigma)|<\operatorname{Re}( \varsigma),~|\operatorname{Im_j}(\tau)|<\operatorname{Re}(\tau),~|\operatorname{Im_j}(\delta)|<\operatorname{Re}(\delta)$
%	\begin{equation}
%	%\begin{split}
%	%\int_{0}^{\zeta}t^{\tau-1}\mathbb{E}^\delta_{ \varsigma,\tau}(\omega t^\upalph%a)dt
%&=\int_{0}^{\zeta_1}t^{\tau_1-1}\mathbb{E}^{\delta_1}_{ \varsigma_1,\tau_1}(\omega_1 t^ \varsigma_1)dt e_1+ \int_{0}^{\zeta_2}t^{\tau_2-1}\mathbb{E}^{\delta_2}_{ \varsigma_2,\tau_2}(\omega_2 t^ \varsigma_2)dte_2\\
%	&=\zeta_1^{\tau_1}\mathbb{E}^{\delta_1}_{ \varsigma_1,\tau_1+1}(\omega_1 \zeta_1^{ \varsigma_1})e_1+\zeta_2^{\tau_2}\mathbb{E}^\delta_{ \varsigma_2,\tau_2+1}(\omega_2 \zeta_2^{ \varsigma_2})e_2\\
%	&=\zeta^{\tau}\mathbb{E}^{\delta_2}_{ \varsigma,\tau+1}(\omega \zeta^ \varsigma).
%		\end{split}
%	\end{equation}
	
%\end{proof}

\noindent Substituting $ \delta=1$ in the above result (\ref{eq:p int1}), we get the following integral relation for the two-parameter bicomplex ML function:

\begin{equation}\label{eq:2pbmlir}
	\int_{\Gamma}t^{\tau-1}\mathbb{E}_{ \varsigma,\tau}(\omega t^ \varsigma)dt=\zeta^{\tau}\mathbb{E}_{ \varsigma,\tau+1}(\omega \zeta^ \varsigma).
	\end{equation}

\section{Integral transforms of Bicomplex Prabhakar Function }
In this section,  bicomplex integral transforms of the bicomplex Prabhakar function are evaluated. The results obtained are instrumental in advancing the theoretical framework of the newly developed bicomplex Prabhakar function. By establishing these integral transforms, we lay the groundwork for further exploration and application in various mathematical and physical contexts. 
%\subsection*{Bicomplex	  Laplace Transform and its Inverse}
%In \cite{akumar2011a}, Kumar et al. presented the bicomplex Laplace transform and its basic properties.
%Kumar et al. \cite{akumar2011a} introduced the bicomplex Laplace transform and its basic properties. %These findings have a lot of potential in the area of signal processing.
%These results can be highly applicable in the field of signal processing.\\
%\begin{definition}[Bicomplex	  Laplace transform]
%	Let $f(t)$ be a bicomplex valued function of exponential order $M\in \mathbb{R}$. Then Laplace
%	transform of $h(t)$ for $t \ge 0$  defined as:
%	\begin{equation}\label{eq:bc lap def}
%	\mathcal{L}\{h(t)\} = \hat{h}(\zeta)= \int_{0}^{\infty}h(t)e^{-\zeta t}dt.
%	\end{equation}
%	Here $\hat{h}(\zeta)$ exist and is convergent for all $\zeta \in  D$
%	where
%		\begin{equation}\label{eq:lap domain}\notag
%	D= \{\zeta:H_\rho(\zeta)~\text{ represents a Right half plane}~ a_0>	M+|a_3|  \},
%	\end{equation}	
 %  and  $\zeta = a_0 + a_1i_1 + a_2i_2 + a_3j.$\\
   
%\noindent   For more details  about bicomplex Laplace transform see \cite{akumar2011a,ragarwal2014a}.
    %Agarwal et al. 
    %\cite{ragarwal2017} define     Bicomplex Mellin transform  was defined by Agarwal et al. \cite{ragarwal2017}  as follows :
%	\begin{definition}[Bicomplex Mellin transform]
%		Let $f(t)$ be a bicomplex-valued continuous function on the interval $(0,\infty)$ with $f(t)=O\left(t^{- \varsigma}\right)$ as $t\to 0^{+}$ and $f(t)=O\left(t^{-\tau}\right)$ as $t\to \infty,$ where $ \varsigma<\tau$. Then bicomplex Mellin transform  \cite{ragarwal2017} of $f(t)$ defined as
%		\begin{equation}
%\mathcal{M}[f(t);\zeta] =\tilde{f}(\zeta)=\int_{0}^{\infty}t^{\zeta-1}f(t)dt,\quad \zeta % \in \Omega,
	%	\end{equation}
%		where $\tilde{f}(\zeta)$ is analytic and %convergent in $\Omega$ defined as
%		\begin{equation}
%		\Omega=\{\zeta\in %\mathbb{T}: \varsigma+\left|\text{Im}_{j}(\zeta)\right|<\text{Re}(\zeta)<\tau-\left|\text{Im}_{j}(\zeta)\right|\}.
%		\label{8n8}
%		\end{equation}
		%where $\text{Im}_{j}(\zeta)$ denotes the imaginary part w.r.t. $j$ unit of a bicomplex number.
	%\end{definition}
	%	In D there are infinite $\zeta$ which have same $H_\rho$ hyperbolic projection because	$a_1$ and $a_2$ are free from restriction.
%end{definition}
\begin{definition}
    The bicomplex Laplace transform (LT) 	of the Riemann-Liouville fractional integral \cite{agarwal2023bicomplex} of the function $f(t)$ (piecewise and of exponential order), is given by
\begin{equation}\label{eq:LT-fracD}
\mathcal{L}\left(  _0D^{-\mu}_tf(t);\zeta \right) =\zeta^{-\mu}\hat{f}(\zeta),~\zeta ,\mu\in\mathbb{T},
%\mathcal{L}\left(  I^{\mu}_{0+}f(t);s\right) =s^{-\mu}\mathcal{F}(s),~s,\mu\in\mathbb{T},
\end{equation}
where $ \hat{f}(\zeta)$ is the LT of $f( t)$, $|\operatorname{Im_j}(\zeta )|<\operatorname{Re}(\zeta ),$$|\operatorname{Im_j}(\mu)|<\operatorname{Re}(\mu).$ 

\end{definition}
\subsection{ Bicomplex Laplace transform of the Prabhakar function}
LT of Prabhakar function is defined by (see, e.g. \cite{gorenflo2014,kilbas2006})
\begin{equation}\label{eq lt 3 para}
\mathcal{L}(t^{\tau-1}E_{ \varsigma,\tau}^\delta(\omega t^{ \varsigma});s)=\int_{0}^{\infty}e^{-st}E_{ \varsigma, \tau}^\delta(\omega  t^{ \varsigma})t^{\tau-1}dt=\frac{s^{ \varsigma\delta-\tau}}{(s^{ \varsigma}-\omega)^\delta},
\end{equation}
where $ \operatorname{Re}(\tau) > 0, ~\operatorname{Re}(s) > 0,~ |\omega s^{- \varsigma}|<1,~\omega \in\mathbb{C}.$\\

\begin{theorem}\label{th:lt two para}
Let $\zeta,~ \varsigma,~\tau,~\delta,~\uplambda \in \mathds{T} $ where $\zeta=\zeta_1e_1+\zeta_2e_2$ bicomplex
LT of the Prabhakar function	is given by 
\begin{equation}\label{eq:bclt 3para}
\mathcal{L}(t^{\tau-1}E_{ \varsigma,\tau}^\delta(\uplambda t^{ \varsigma}))(\zeta)=\int_{0}^{\infty}e^{-\zeta t}E_{ \varsigma, \tau}^\delta(\uplambda t^{ \varsigma})t^{\tau-1}dt=\frac{\zeta^{ \varsigma\delta-\tau}}{(\zeta^{ \varsigma}-\uplambda)^\delta},
\end{equation}
where	
$|\operatorname{Im_j}(\tau)|<\operatorname{Re}(\tau),~|\operatorname{Im_j}(\zeta)|<\operatorname{Re}(\zeta),~|\uplambda \zeta^{- \varsigma}|_j\prec1.$\\
	
\end{theorem}
\begin{proof}
Making use of idempotent representations of $\zeta,~ \varsigma,~\tau,~\lambda \in \mathds{T} $ %where $\zeta=\zeta_1e_1+\zeta_2e_2$. 
and with the help of relation 
 (\ref{eq lt 3 para}) and the result (\ref{eq: bc ml three para}), we have
	\begin{equation}\label{eq:bc lap 2mla}
	\begin{split}
	\mathcal{L}(t^{\tau-1}E_{ \varsigma,\tau}^\delta(\uplambda t^{ \varsigma});\zeta)&=\mathcal{L}(t^{\tau_1-1}E_{ \varsigma_1,\tau_1}^{\delta_1}(\uplambda_1 t^{ \varsigma_1});\zeta_1)e_1+\mathcal{L}(t^{\tau_2-1}E_{ \varsigma_2,\tau_2}^{\delta_2}(\uplambda_2 t^{ \varsigma_2});\zeta_2)e_2\\
	&=\int_{0}^{\infty}e^{-\zeta_1 t}E_{ \varsigma_1, \tau_1}^{\delta_1}(\uplambda_1 t^{ \varsigma_1})t^{\tau_1-1}dte_1+\int_{0}^{\infty}e^{-\zeta_2 t}E_{ \varsigma_2, \tau_2}^{\delta_2}(\uplambda_2 t^{ \varsigma_2})t^{\tau_2-1}dte_2\\
	&=\frac{\zeta_1^{ \varsigma_1\delta_1-\tau_1}}{(\zeta_1^{ \varsigma_1}-\uplambda_1)^{\delta_1}}e_1+\frac{\zeta_2^{ \varsigma_2\delta_2-\tau_2}}{(\zeta_2^{ \varsigma_2}-\uplambda_2)^{\delta_2}}e_2,~~~|\uplambda_1 \zeta_1^{- \varsigma_1}|<1,~|\uplambda_2 \zeta_2^{- \varsigma_2}|<1\\
	&=\frac{\zeta^{ \varsigma\delta-\tau}}{(\zeta^{ \varsigma}-\uplambda)^\delta}.
	\end{split}
	\end{equation}
	Here,
	\begin{equation}
	%\begin{split}
	|\uplambda \zeta^{- \varsigma}|_j=|\uplambda_1 \zeta_1^{- \varsigma_1}|e_1+|\uplambda_2 \zeta_2^{- \varsigma_2}|e_2
	\prec1.e_1+1.e_2
	=1.
	%\end{split}
	\end{equation}	
	
\end{proof}
\subsection{Bicomplex Mellin transform of Prabhakar function}
The Mellin transform of the Prabhakar function is given by \cite{gorenflo2014}
\begin{equation}\label{eq mt 3 para}
\mathcal{M}(E_{ \varsigma,\tau}^\delta( -\omega t;s))=\int_{0}^{\infty}t^{s-1}E_{ \varsigma, \tau}^\delta(-\omega t)dt=\frac{\Gamma(s)\Gamma(\delta-s)}{\Gamma(\delta)\Gamma(\tau- \varsigma s)}\omega^{-s},
\end{equation}
where $ \varsigma\in \mathbb{R}_+, ~\omega,\tau,\delta\in \mathbb{C},$  {$\tau\neq 0$},~ $\operatorname{Re}(\delta) > 0.$\\

\noindent Here, we evaluate the bicomplex Mellin transform of Prabhakar function in the bicomplex space.
\begin{theorem}\label{th:mt three para}
	Let $ \varsigma\in \mathbb{R}_+,~\zeta,~\uplambda ,~\tau,~\delta\in \mathds{T} $ where $\zeta=\zeta_1e_1+\zeta_2e_2, ~\uplambda =\uplambda _1e_1+\uplambda _2e_2,~\tau=\tau_1e_1+\tau_2e_2 .$ Bicomplex
MT of Prabhakar function	is given by 
\begin{equation}\label{eq mt 3 para 1}
\mathcal{M}(E_{ \varsigma,\tau}^\delta( -\uplambda t))(\zeta)=\int_{0}^{\infty}t^{\zeta-1}E_{ \varsigma, \tau}^\delta(-\uplambda t)dt=\frac{\Gamma(\zeta)\Gamma(\delta-\zeta)}{\Gamma(\delta)\Gamma(\tau- \varsigma \zeta)}\uplambda ^{-\zeta},
\end{equation}
where	
	$|\operatorname{Im_j}(\delta)|<\operatorname{Re}(\delta), ~ \operatorname{Re}(\tau)\neq |\operatorname{Im_j}(\tau)|~\text{and}~ \operatorname{Im_{i_1}}(\tau)\neq |\operatorname{Im_{i_2}}(\tau)|.$\\
	
\end{theorem}
\begin{proof}
	Let $\zeta,~ \varsigma,~\tau,~\uplambda  \in \mathds{T} $ where $\zeta=\zeta_1e_1+\zeta_2e_2$
By the  relation \eqref{eq mt 3 para} and the result \eqref{eq: bc ml three para}, using the idempotent representation for the bicomplex Prabhakar function and simplifying the existence conditions, we obtain
\begin{equation}\label{eq:bc lap 2ml}
\begin{split}
	\mathcal{M}(E_{ \varsigma,\tau}^\delta( -\uplambda t))(\zeta)
&=	\int_{0}^{\infty}t^{\zeta_1-1}E_{ \varsigma, \tau_1}^{\delta_1}(-\uplambda _1t)dt~e_1+\int_{0}^{\infty}t^{\zeta_2-1}E_{ \varsigma, \tau_2}^{\delta_2}(-\uplambda _2t)dt~e_2\\
&=\frac{\Gamma(\zeta_1)\Gamma(\delta_1-\zeta_1)}{\Gamma(\delta_1)\Gamma(\tau_1- \varsigma \zeta_1)}\uplambda _1^{-\zeta_1}~e_1+\frac{\Gamma(\zeta_2)\Gamma(\delta_2-\zeta_2)}{\Gamma(\delta_2)\Gamma(\tau_2- \varsigma \zeta_2)}\uplambda _2^{-\zeta_2}~e_2\\
&=\frac{\Gamma(\zeta)\Gamma(\delta-\zeta)}{\Gamma(\delta)\Gamma(\tau- \varsigma \zeta)}\uplambda ^{-\zeta},
	\end{split}
	\end{equation}
	%where $|\operatorname{Im_j}(\delta)|<\operatorname{Re}%(\delta).$
	where $|\operatorname{Im_j}(\delta)|<\operatorname{Re}(\delta)$ and
$\operatorname{Re}(\tau)\neq |\operatorname{Im_j}(\tau)|~ \operatorname{Im_{i_1}}(\tau)\neq |\operatorname{Im_{i_2}}(\tau)|.$\\
	Here
	\begin{equation}
	\begin{split}
	\tau &= p_0+i_1p_1+ i_2p_2+i_1i_2p_3
	%&=(a_0+i_1a_1) +i_2( a_2+i_1a_3)\\
	=\tau_1e_1+\tau_2e_2,
	\end{split}
	\end{equation}
	where $\tau_1= (p_0+p_3) +i_1(p_1-p_2)$ and $\tau_2= (p_0-p_3) +i_1(p_1+p_2).$
	%and $ ~|Im( \varsigma_2)|<Re( \varsigma_1).$ 
	
\noindent	Since
$\tau_1\ne0,~\tau_2 \ne0.$
\begin{eqnarray}
&\Rightarrow & p_0+p_3 \neq 0, ~ p_0-p_3\neq 0~\text{and}~ p_1-p_2 \neq 0 ,~ p_1+p_2\neq 0 \notag\\
&	\Rightarrow& ~ p_0\neq |p_3|~\text{and}~ p_1\ne |p_2|\notag\\
%&\Rightarrow&|\operatorname{Im_j}( \varsigma)|<\operatorname{Re}( \varsigma)\\
&\Rightarrow& \operatorname{Re}(\tau)\neq |\operatorname{Im_j}(\tau)|~\text{and}~ \operatorname{Im_{i_1}}(\tau)\neq |\operatorname{Im_{i_2}}(\tau)|.
\end{eqnarray}


\end{proof}

\section{Application of Bicomplex Prabhakar Function}
	

\noindent The unified fractional kinetic equation is versatile and finds applications in various fields. It models anomalous diffusion, stochastic processes, and viscoelastic materials. It's also used in biological systems, astrophysics, fractional boundary-value problems, and financial mathematics, helping to understand complex behaviors and solve practical problems \cite{saichev1997}.\\
In this section is  the solution of
the unified fractional kinetic equation has been derived in the terms of bicomplex Prabhakar function. %This method is based on the application of bicomplex Laplace transform.
\begin{theorem}
	If $|\operatorname{Im_j}(\upnu_i)|<\operatorname{Re}(\upnu_i),~ a_i (a_{i_1}> |a_{i_4}|,~i\in\mathbb{N})$ are hyperbolic numbers   and function   $f\in L(\mathbb{R}^+)$ is bicomplex valued integrable function defined for $t\in\mathbb{R}^+$. Then the equation
	\begin{equation}\label{eq: kinetic1}
	N(t)-N_0f(t)=-\sum_{i=1}^{n}{a_i}~  _0D^{-\upnu_i}_tN(t),
	\end{equation}
	is solvable and its particular solution is given by
	\begin{equation}\label{eq: kinetic2}
	\begin{split}
		N(t)&= N_0\sum_{l=0}^{\infty}(-1)^l \sum_{s_1+\dots+s_{n-1}=l}\frac{(l)!}{(s_1)!\dots(s_{n-1})!}\left( \prod_{\upmu=1}^{n-1}(a_{\upmu+1})^{s_\upmu})\right) \\
	&\times \int_{0}^{t}f(v)(t-v)^{\sum_{\upmu=1}^{n-1}\upnu_{\upmu +1}-1}\mathbb{E}_{\upnu_1,\sum_{\upmu=1}^{n-1}\upnu_{\upmu +1}}^{l+1} \left( -a_1(t-v)^{\upnu_1}\right) dv,
	\end{split}
	\end{equation}
where the summation in (\ref{eq: kinetic2}) is taken over all non-negative integers $s_1,\dots ,s_n$ such that $s_1+\dots +s_{n-1}=l$ and provided that the series and integral in (\ref{eq: kinetic2}) are convergent. 
\end{theorem}
\begin{proof}
By taking LT of (\ref{eq: kinetic1}) and using the result (\ref{eq:LT-fracD}) we get
\begin{equation}
\mathcal{L}\left( 	N(t)-N_0f(t)\right) =-\sum_{i=1}^{n}{a_i}~ \mathcal{L}\left(  _0D^{-\upnu_i}_tN(t)\right) 
\end{equation}	
\begin{equation}
\Rightarrow \hat { N}(\zeta) -N_0\hat{f}(\zeta)=-\sum_{i=1}^{n}{a_i}~ \zeta^{-\upnu_i}\hat { N}(\zeta)
\end{equation}
\begin{equation}\label{eq:kin}
\begin{split}
\Rightarrow \hat { N}(\zeta)&=\frac{N_0\hat{f}(\zeta)}{1+a_1 \zeta^{-\upnu_1}+\dots+a_n\zeta^{-\upnu_n}}\\
&=N_0\hat{f}(\zeta)(-1)^l\frac{\left( \sum_{j=1}^{n-1}a_{j+1}\zeta^{-\upnu_{j+1}}\right)^l }{(1+a_1\zeta^{-\upnu_1})^{l+1}},~\left| \frac{\sum_{j=2}^{n}a_{j}\zeta^{-\upnu_{j}}}{1+a_1\zeta^{-\upnu_{1}}}\right| <1.
\end{split}
\end{equation}
By using the identity (see, e.g.\cite{saxena2010})
\begin{equation}
(x_1+\dots+x_m)^l=\sum_{s_1+\dots+s_{n}=l} \frac{(l)!}{(s_1)!\dots (s_n)!}\prod_{\upmu=1}^{n}x_\upmu^{s_\upmu},
\end{equation}
where the summation  is taken over all non-negative integers $s_1,\dots, s_n$ such that $s_1+\dots+ s_{n}=l$ then for $|a_1\zeta^{-\upnu_{1}}|\prec1$ from equation (\ref{eq:kin})  we get
\begin{equation}
\hat { N}(\zeta)=N_0\hat{f}(\zeta)(-1)^l \sum_{\substack{s_1+\dots+s_{n_1}=l\\s_1>\dots s_{n_1}}} \frac{(l)!}{(s_1)!\dots (s_{n-1})!} \frac{\left( \prod_{\upmu=1}^{n-1}(a_{\upmu+1})^{s_\upmu}\right)\zeta^{-\sum_{\upmu=1}^{n-1}\upnu_{\upmu+1}}}{(1+a_1\zeta^{-\upnu_{1}})^{l+1}}. 
\end{equation}
Taking the inverse LT of the above equation and by making use of the result (\ref{eq:bclt 3para}) and using the convolution theorem %of the Laplace transform, 
the result (\ref{eq: kinetic2}) is obtained (see, e.g. \cite{saxena2010}).\\
	Also
\begin{equation}\label{eq:an1}
\begin{split}
a_i&=a_{i_1} +i_1a_{i_2}+i_2a_{i_3}+j a_{i_4}\\
&=(a_{i_1} +i_1a_{i_2})+i_2(a_{i_3}+i_1a_{i_4})\\
&=b_i e_1+c_i e_2.
\end{split}
\end{equation}
	Here,  $b_i= (a_{i_1}+a_{i_4}) +i_1(a_{i_2}-a_{i_3})$ and $c_i= (a_{i_1}-a_{i_4}) +i_1(a_{i_2}+a_{i_3}).$\\
	Since,
\begin{eqnarray}
&\Rightarrow&b_i>0 \text{ and} ~ c_i=0.\notag \\
&\Rightarrow& a_{i_1}+a_{i_4}>0, ~a_{i_2}-a_{i_3}=0~\text{and}~ a_{i_1}-a_{i_4}>0 ,~a_{i_2}+a_{i_3}=0. \notag \\
&	\Rightarrow& a_{i_1}> |a_{i_4}|~\text{and}~ a_{i_2}=a_{i_3}=0.\\
&	\Rightarrow &	a_i=(a_{i_1} +a_{i_4})e_1+(a_{i_1} -a_{i_4})e_2=a_{i_1} +ja_{i_4}.\label{eq :an equi}	
\end{eqnarray}	
%$\Rrightarrow a_{i_1}+a_{i_4}>0, ~a_{i_2}-a_{i_3}=0 $ aid $a_{i_1}-a_{i_4}>0 ,~a_{i_2}+a_{i_3}=0. $\\
%$\Rrightarrow a_{i_1}> |a_{i_4}|$ aid $a_{n_2}=a_{n_3}=0.$\\
%$\Rrightarrow 	a_n=(a_{n_1} +a_{n_4})e_1+(a_{n_1} -a_{n_4})e_2=a_{n_1} +ja_{n_4}.$\\
Hence, $a_i $ is a hyperbolic number such that $a_{i_1}> |a_{i_4}|.$ \\


\end{proof}

\noindent For $f(t)=t^{\tau -1}\mathbb{E}_{ \varsigma,\tau}^\delta[-(a t)^{ \varsigma}]$, as a particular case, we obtain following  special result involving the bicomplex Prabhakar function: 

\begin{theorem}
	Let $ \varsigma ,~\tau\in \mathbb{T},~|\operatorname{Im_j}( \varsigma)|<\operatorname{Re}( \varsigma),~ |\operatorname{Im_j}(\tau)|<\operatorname{Re}(\tau),~i\in\mathbb{N}, ~a (u_{1}> |u_{4}|)$ is hyperbolic number  and $\mathbb{E}_{ \varsigma,\tau}^\delta(\zeta)$ is bicomplex Prabhakar function then the equation
	\begin{eqnarray}\label{ke-ml}
	N(t)-N_0 t^{\tau -1}\mathbb{E}_{ \varsigma,\tau}^\delta[-(a t)^{ \varsigma}]=-\sum_{r=1}^{n}\binom{n}{r} {a}^{r \varsigma}~  _0D^{-r \varsigma}_t N(t),~n\in \mathbb{N},
	\end{eqnarray}
	holds the relation
	\begin{equation}\label{eq:app}
	N(t) = N_0t^{\tau -1}\mathbb{E}_{ \varsigma,\tau}^{\delta+n}[-(a t)^{ \varsigma}],~n\in \mathbb{N}.
	\end{equation}
\end{theorem}
\begin{proof} Taking Laplace transform of the equation \eqref{ke-ml} and using the results \eqref{eq:LT-fracD} and \eqref{eq:bclt 3para} therein, we obtain
\begin{equation}
\begin{split}
\hat { N}(\zeta)-N_0 \mathcal{L}\left( t^{\tau -1}\mathbb{E}_{ \varsigma,\tau}^\delta[-(a t)^{ \varsigma}];\zeta\right) &=-\sum_{r=1}^{n}\binom{n}{r} {a}^{r \varsigma}~ \zeta^{-r \varsigma} \hat { N}(\zeta).\\
\Rightarrow \hat { N}(\zeta)\left[ 1+\sum_{r=1}^{n}\binom{n}{r} {a}^{r \varsigma}~ \zeta^{-r \varsigma}\right]& =N_0 \frac{\zeta^{ \varsigma\delta-\tau}}{(\zeta^{ \varsigma}+a^ \varsigma)^\delta}.\\
\Rightarrow \hat { N}(\zeta)\left[ 1+\sum_{r=1}^{n}\binom{n}{r} (\zeta/a)^{-r \varsigma}\right]& =N_0 \frac{\zeta^{ \varsigma\delta-\tau}}{(\zeta^{ \varsigma}+a^ \varsigma)^\delta}.\\
\Rightarrow \hat { N}(\zeta)\left[ 1+\sum_{r=1}^{n}\binom{n}{r}[ (\zeta/a)^{- \varsigma}]^r\right]& =N_0 \frac{\zeta^{ \varsigma\delta-\tau}}{(\zeta^{ \varsigma}+a^ \varsigma)^\delta}.\\
\Rightarrow \hat { N}(\zeta)\left[ 1+ (\zeta/a)^{- \varsigma}\right]^n& =N_0 \frac{\zeta^{ \varsigma\delta-\tau}}{(\zeta^{ \varsigma}+a^ \varsigma)^\delta}.\\
\Rightarrow \hat { N}(\zeta)& =N_0 \frac{\frac{\zeta^{ \varsigma\delta-\tau}}{(\zeta^{ \varsigma}+a^ \varsigma)^\delta}}{\left[ 1+ (\zeta/a)^{- \varsigma}\right]^n}.\\
\Rightarrow \hat { N}(\zeta)& =N_0 \frac{\frac{\zeta^{ \varsigma\delta-\tau}}{(\zeta^{ \varsigma}+a^ \varsigma)^\delta}}{\left[ \frac{\zeta^ \varsigma+ a^{ \varsigma}}{\zeta^ \varsigma}\right]^n}.\\
\Rightarrow \hat { N}(\zeta)& =N_0 \frac{\zeta^{ \varsigma\delta-\tau}}{(\zeta^{ \varsigma}+a^ \varsigma)^\delta}{\left[ \frac{\zeta^ \varsigma}{\zeta^ \varsigma+ a^{ \varsigma}}\right]^n}.\\
%\Rightarrow \hat { N}(\zeta)& =N_0 \frac{\zeta^{ \varsigma(\delta+n)-\tau}}{(\zeta^{ \varsigma}+a^ \varsigma)^{\delta+n}}\label{eq:invlap}\\
\end{split}
\end{equation}
Hence
\begin{equation}\label{eq:invlap}
    \hat { N}(\zeta) =N_0 \frac{\zeta^{ \varsigma(\delta+n)-\tau}}{(\zeta^{ \varsigma}+a^ \varsigma)^{\delta+n}}.
\end{equation}
By taking inverse Laplace transform of \eqref{eq:invlap} we get
\begin{equation}\label{eq:app1}
	N(t) = N_0t^{\tau -1}\mathbb{E}_{ \varsigma,\tau}^{\delta+n}[-(a t)^{ \varsigma}],~n\in \mathbb{N}.
	\end{equation}
%By  breaking up into idempotent components taking inverse Laplace transform of \eqref{eq:inv} 
%by  using the  following formula  (see, e.g. \cite{saxena2010}) and  using the Corollary 2.2 \cite{saxena2010}. Combining the idempotent omponents to obtain the solution in the bicomplex domain we get %result \eqref{eq:app}.
%\begin{equation}
%\mathcal{L}^{-1}[s^{-\tau}(1-as^{- \varsigma})^{-\delta});t]=t^{\tau-1}\mathbb{E}_{ \varsigma,\tau}^\delta (at^ \varsigma)
%\end{equation}


\end{proof}
    

\begin{remark}
Since
\begin{equation}\label{eq: an}
\begin{split}
a&=u_{1} +i_1u_{2}+i_2u_{3}+j u_{4}\\
&=(u_{1} +i_1u_{2})+i_2(u_{3}+i_1u_{4})\\
&=b e_1+c e_2.
\end{split}
\end{equation}
Here,  $b= (u_{1}+u_{4}) +i_1(u_{2}-u_{3})$ and $c= (u_{1}-u_{4}) +i_1(u_{2}+u_{3}).$\\
Since,
\begin{eqnarray}
&\Rightarrow&b>0 \text{ and} ~ c=0.\notag \\
&\Rightarrow& u_{1}+u_{4}>0, ~u_{2}-u_{3}=0~\text{and}~ u_{1}-u_{4}>0 ,~u_{2}+u_{3}=0. \notag \\
&	\Rightarrow& u_{1}> |u_{4}|~\text{and}~ u_{2}=u_{3}=0.\\
&	\Rightarrow &	a=(u_{1} +u_{4})e_1+(u_{1} -u_{4})e_2=u_{1} +ju_{4}.\label{eq :an equi1}	
\end{eqnarray}	
%$\Rrightarrow a_{1}+a_{4}>0, ~a_{2}-a_{3}=0 $ aid $a_{1}-a_{4}>0 ,~a_{2}+a_{3}=0. $\\
%$\Rrightarrow a_{1}> |a_{4}|$ aid $a_{n_2}=a_{n_3}=0.$\\
%$\Rrightarrow 	a_n=(a_{n_1} +a_{n_4})e_1+(a_{n_1} -a_{n_4})e_2=a_{n_1} +ja_{n_4}.$\\
Hence, $a $ is a hyperbolic number such that $u_{1}> |u_{4}|.$ \\
\end{remark}
%\end{proof}
\section{Conclusion}
This paper defines the Prabhakar function and its properties in bicomplex space, extending from its complex counterpart. It derives various properties, special cases, recurrence relations, integral representations, and differential relations. The application of the bicomplex LT of the bicomplex Prabhakar function is illustrated to solve fractional kinetic equations. Further research can explore more results on the Prabhakar function in bicomplex space, as this area remains largely unexplored. Additionally, fractional calculus involving the Prabhakar function in the kernel presents a promising topic for future study.
% The Prabhakar integral, Prabhakar derivative, and regularised Prabhakar derivative in bicomplex space have been discussed in this paper.
\\
	
\textbf{Funding declaration:} {This research received no external funding.}\\
	
	
%\dataavailability{\hl{No} new data were created or analyzed in this study. Data sharing is not applicable to this article.} 
	
	
%\acknowledgments{ \hl{The} %MDPI: Please ensure that all individuals included in this section have consented to the acknowledgement.

	
\textbf{Conflicts of interest:} The authors declare no conflicts of {interest}.\\
%No, I declare that the authors have no competing interests as defined by Springer, or other interests that might be perceived to influence the results and/or discussion reported in this paper.

\bibliographystyle{apalike}%unsrt
\bibliography{lref}


\end{document}
%https://www.mdpi.com/2504-3110/8/9/508
Extensions of Bicomplex Hypergeometric Functions and Riemann–Liouville Fractional Calculus in Bicomplex Numbers
